\documentclass[9pt]{elife}
\usepackage{amsthm}
\usepackage{amsmath,amssymb,mathtools}
\usepackage{siunitx}
\usepackage{graphicx}
\usepackage{booktabs}
\usepackage{hyperref}
\usepackage{tcolorbox}
\usepackage{bm}

\title{MacroIR: A Bayesian framework for analytically exact interval likelihoods in time-averaged biological data}

\author[1*]{Luciano Moffatt}
\affil[1]{Instituto de Qu\'imica F\'isica de los Materiales, Medio Ambiente y Energ\'ia (INQUIMAE), CONICET; Facultad de Ciencias Exactas y Naturales, Universidad de Buenos Aires, Ciudad de Buenos Aires, C1428EHA, Argentina}
\corr{lmoffatt@qi.fcen.uba.ar}

\begin{document}
	\maketitle
	\begin{abstract}
		Time-averaged observations are pervasive in experimental biology, yet most inference frameworks assume instantaneous sampling and thus misestimate systems whose dynamics evolve faster than the acquisition window. 
		\textbf{MacroIR} (Macroscopic Interval Recursive) formulates a fully Bayesian likelihood for interval-averaged signals generated by continuous-time Markov processes. 
		Exploiting the Markov property, the method treats each observation window as a bridge between an initial and a final state. 
		Given the prior at the start and the kinetic parameters, this pair suffices to predict the probability of all trajectories within the interval. 
		The observed average depends jointly on both states, enabling a Bayesian update over the Cartesian distribution of start–end pairs using the analytical mean and variance of the observable. 
		Marginalizing over the initial state yields the posterior for the next interval, completing a deterministic recursion without Monte Carlo sampling. 
		In the linear–Gaussian limit, MacroIR reduces exactly to the Kalman filter, clarifying the latter as a special case of Bayesian inference on Markov processes. 
		This formulation unifies stochastic kinetics, time-averaged observables, and Bayesian evidence evaluation, providing unbiased parameter and state inference across molecular and cellular datasets.
	\end{abstract}
	\begin{abstract}
		Time-averaged measurements dominate experimental biology, yet most inference methods assume instantaneous observations, yielding biased estimates when intrinsic dynamics outpace acquisition. 
		\textbf{MacroIR} (Macroscopic Interval Recursive) provides a fully Bayesian formulation of the likelihood for interval-averaged signals generated by continuous-time Markov processes. 
		By conditioning on both the start and end states of each sampling window, MacroIR computes the interval mean and variance in closed form and recursively updates the posterior probability distribution over hidden states. 
		Under a Gaussian-moment approximation, these updates reduce to the familiar Kalman recursion, but the Bayesian derivation remains valid beyond the linear–Gaussian regime. 
		The framework unifies ensemble kinetics, time-averaged observables, and Bayesian evidence computation. 
		Applications to simulated and experimental data demonstrate that MacroIR eliminates structural bias in time-averaged biological recordings and enables reliable inference of underlying dynamics across molecular, cellular, and behavioral timescales.
	\end{abstract}
	
	\section*{Introduction}
	
	Time-averaged measurements are ubiquitous in experimental biology—electrophysiological amplifiers integrate currents over microseconds, fluorescence detectors average photon flux over milliseconds, and population assays aggregate signals over hours \citep{Hille2001, Denk1990, Elowitz2002}. 
	Yet most inference methods treat such data as instantaneous samples. 
	When intrinsic dynamics evolve on comparable timescales, this mismatch between physical averaging and model assumptions distorts likelihoods and biases inferred parameters \citep{ColquhounHawkes1982, Qin2000, Munsky2009, Moffatt2007}.
	
	Continuous-time Markov models (CTMCs) remain the standard framework for describing stochastic molecular and cellular dynamics. 
	Over recent decades, inference methods have evolved from rate-equation and hidden-Markov formulations to recursive Bayesian estimators that continuously update knowledge of the system’s state \citep{Csanady2006, Siekmann2011, Bronson2013}.  
	Münch \emph{et al.} (2022) recently introduced a generalized Bayesian filter for ion-channel kinetics that unifies current and fluorescence data and explicitly integrates mean dynamics over finite sampling windows \citep{Munch2022}.  
	However, their approach retains an implicit Gaussian approximation for the posterior distribution and assumes instantaneous observation within each bin, leading to bias when acquisition windows approach the system’s dwell times.
	
	A direct precursor of this work was the \emph{macroscopic recursive} (Macro R) algorithm introduced by Moffatt (2007) for the Bayesian estimation of kinetic rates from macroscopic current fluctuations.  
	Macro R formulated the full recursive likelihood for ensemble recordings, approximating the hidden state distribution by a Gaussian and updating its moments after each measurement.  
	The resulting recursion was mathematically identical to a Kalman filter but conceptually derived from first Bayesian principles.  
	The present framework, \textbf{MacroIR} (Macroscopic Interval Recursive), extends that Bayesian theory to \textbf{time-averaged observations}, deriving analytically exact interval likelihoods and recursive updates for the posterior distribution.  
	This extension restores consistency between the physical measurement process and the probabilistic model, allowing unbiased inference and rigorous model comparison via Bayesian evidence.
	
	\section*{Theory}
	
	\subsection*{From the macroscopic recursive (Macro R) to the interval-recursive (Macro IR) formulation}
	
	In the Macro R approach \citep{Moffatt2007}, the ensemble of channel occupancies is represented by a probability distribution $p(\mathbf{s},t)$ over discrete states that evolves under a rate matrix $Q$.  
	After each instantaneous measurement $y_t$, Bayes’ rule updates the posterior:
	\[
	p^{\text{post}}(\mathbf{s}|y_t) \propto p(y_t|\mathbf{s})\,p^{\text{prior}}(\mathbf{s}).
	\]
	When the ensemble contains many channels, $p(\mathbf{s})$ can be approximated by a multivariate normal characterized by mean vector $\bm{\mu}$ and covariance matrix $\Sigma$.  
	The propagation of these moments between observations yields what is commonly known as the Kalman recursion but here interpreted as the \emph{Gaussian-moment limit} of Bayesian inference.
	
	MacroIR generalizes this logic to time-averaged observations.  
	Instead of observing $y_t$ at a single point, detectors measure the mean observable over an acquisition window $[0,t]$:
	\[
	\bar y_{0\to t} = \frac{1}{t}\int_0^t \gamma(x_s)\,ds.
	\]
	The likelihood must therefore condition on both the start and end states of the hidden Markov process within each interval.
	
	\subsection*{Boundary-conditioned Bayesian formulation (meta-state representation)}
	
	Let $x_t$ denote a CTMC with generator $Q$ and propagator $P(t)=e^{Qt}$.  
	For a linear observable $\bm{\gamma}\in\mathbb{R}^k$, the conditional mean of the interval average, given start state $i$ and end state $j$, is
	\begin{equation}
		\gamma_{i\to j}(t)
		= \frac{1}{P_{i\to j}(t)}
		\sum_{n_1,n_2}
		V_{i n_1}\,(V^{-1}\Gamma V)_{n_1 n_2}\,V^{-1}_{n_2 j}\,
		E_2(\lambda_{n_1} t, \lambda_{n_2} t),
	\end{equation}
	where $Q=V\Lambda V^{-1}$, $\Gamma=\mathrm{diag}(\bm{\gamma})$, and
	\[
	E_2(x,y)=\frac{e^x-e^y}{x-y}.
	\]
	An analogous expression for the variance involves a three-exponential kernel $E_3(x,y,z)$ derived from the spectral decomposition of $Q$:
	\begin{align}
		\mathrm{Var}(\Gamma_{i\to j}) 
		= E[\Gamma^2_{i\to j}] - (E[\Gamma_{i\to j}])^2,
	\end{align}
	\begin{align}
		E[\Gamma^2_{i\to j}] 
		= \frac{1}{P_{i\to j}(t)}
		\sum_{k_1,k_2,n_1,n_2,n_3}
		V_{i n_1}V^{-1}_{n_1 k_1}\gamma_{k_1}
		V_{k_1 n_2}V^{-1}_{n_2 k_2}\gamma_{k_2}
		V_{k_2 n_3}V^{-1}_{n_3 j}
		E_3(\lambda_{n_1}t,\lambda_{n_2}t,\lambda_{n_3}t).
	\end{align}
	
	\begin{tcolorbox}[colback=gray!5,colframe=gray!60,title=Boundary-conditioned two-point states (“meta-states”)]
		Define the ordered pair $(i,j)$ representing start and end states over duration $t$.  
		The $k^2$ possible pairs form a space whose generator is the Kronecker sum
		$Q\oplus Q = Q\otimes I + I\otimes Q$.  
		Because the propagator factorizes as $e^{(Q\oplus Q)t}=e^{Qt}\otimes e^{Qt}$, the integrals for the interval mean and variance reduce to bilinear forms in the original $k$-dimensional state space—no explicit $k^2$ object is needed.
	\end{tcolorbox}
	
	This boundary-conditioned Bayesian formulation provides the exact likelihood of observing the interval-averaged measurement $\bar y_{0\to t}$ given $(i,j)$.  
	In the limit $t\!\to\!0$, $\gamma_{i\to j}(t)\!\to\!\gamma_i$ and the formulation collapses to the instantaneous Macro R likelihood.
	
	\subsection*{Gaussian-moment Bayesian recursion}
	
	For an ensemble of $N_{\mathrm{ch}}$ independent channels, the joint occupancy distribution can be approximated by a multivariate normal with mean $\bm{\mu}$ and covariance $\Sigma$.  
	The first two moments of the meta-state Bayesian recursion yield a tractable Gaussian-moment update for the posterior distribution.
	
	Given a prior $(\bm{\mu}^{\mathrm{prior}}_0,\Sigma^{\mathrm{prior}}_0)$ and an observed interval-averaged current $y^{\mathrm{obs}}_{0\to t}$, the predicted mean and variance of the observable are
	\begin{align}
		y^{\mathrm{pred}}_{0\to t} &= N_{\mathrm{ch}}\,(\bm{\mu}^{\mathrm{prior}}_0\!\cdot\!\bm{\gamma}_0),\\
		\sigma^2_{y^{\mathrm{pred}}_{0\to t}} &= 
		\frac{\epsilon^2}{t} + \nu^2_{\mathrm{white}} + \nu^2_{\mathrm{pink}}
		+ N_{\mathrm{ch}}\!\left[\bm{\gamma}_0^\top(\Sigma^{\mathrm{prior}}_0-\mathrm{diag}\,\bm{\mu}^{\mathrm{prior}}_0)\bm{\gamma}_0 
		+ \bm{\mu}^{\mathrm{prior}}_0\!\cdot\!\mathbf{E}\right],
	\end{align}
	where $\bm{\gamma}_0=\{\gamma_{i\to j}(t)\}$ are the interval-averaged conductances derived above.  
	
	The posterior moments follow directly from Bayes’ rule under Gaussian closure:
	\begin{align}
		\bm{\mu}^{\mathrm{post}}_{t} &=
		\bm{\mu}^{\mathrm{prior}}_0 P(t)
		+ \frac{y^{\mathrm{obs}}_{0\to t}-y^{\mathrm{pred}}_{0\to t}}{\sigma^2_{y^{\mathrm{pred}}_{0\to t}}}\,K^\mu_{0\to t},\\
		\Sigma^{\mathrm{post}}_{t} &=
		\Sigma^{\mathrm{prior}}_{t}
		- \frac{1}{\sigma^2_{y^{\mathrm{pred}}_{0\to t}}}
		(K^\mu_{0\to t})^\top K^\mu_{0\to t},
	\end{align}
	with the optimal linear estimator
	\[
	K^\mu_{0\to t} = \frac{\mathrm{Cov}(x_t,\bar y_{0\to t})}{\mathrm{Var}(\bar y_{0\to t})}.
	\]
	When the system is linear and noise is Gaussian, these equations coincide with the classical Kalman filter; otherwise, they remain valid as an approximate Bayesian recursion.
	
	\section*{Methods Old}\label{sec:methods}
	
	In this section, we detail our Bayesian analysis framework, including the experimental data used to evaluate our approach, the kinetic schemes developed for prediction, the likelihood function quantifying predictive accuracy, and the algorithm for computing Bayesian evidence to compare alternative kinetic models and infer posterior parameter distributions. In addition, we describe how the conformational ensemble of the closed P2X receptor—docked with an ATP molecule—was sampled using Molecular Dynamics simulations.
	
	\subsection{Bayesian Analysis}
	\subsubsection{Why a Bayesian Approach?}
	A Bayesian framework is ideally suited for analyzing macroscopic currents because it directly addresses the challenge of temporal autocorrelation—a consequence of successive observations on the same finite population of channels in a macropatch. Both the MacroR and MacroIR algorithms exploit this autocorrelation to extract detailed kinetic information by continuously updating our belief about the probability distribution over the possible kinetic states.
	
	At each stage, these algorithms iteratively refine the posterior probability distribution—approximated as a normal distribution—to yield a dynamic, high-resolution estimate of the kinetic states. Notably, this process also produces a correlation matrix among the states of different channels. It is important to emphasize that this correlation does not imply a causal influence between neighboring channels; rather, it arises because the aggregate current reflects the sum of individual channel contributions, thereby coupling the information on channel states in the observer’s inference. In this sense, the phenomenon is reminiscent of quantum entanglement, where correlations become apparent only after measurement.
	\subsubsection{Objective of the Bayesian Analysis}
	The primary objective of our Bayesian analysis is to extract mechanistic insights from the stochastic fluctuations observed in macrocurrents. We achieve this through two complementary strategies:
	
	\begin{enumerate}[label=\textbf{\arabic*.}]
		\item \textbf{Comparison of Alternative Molecular Mechanisms and Kinetic Schemes.}  
		We translate each proposed molecular mechanism into a corresponding kinetic scheme. In addition, we include previously proposed kinetic schemes where the underlying molecular mechanism remains ambiguous. Bayesian evidence for each alternative is computed using our MCMC framework, and the kinetic schemes that yield the highest evidence are interpreted as having the strongest support from the experimental data. It is important to note that this analysis does not prove that a particular molecular mechanism is occurring; rather, it facilitates selection among a set of plausible alternatives. Definitive validation is achievable only when the set of alternatives encompasses all logically possible options.
		
		\item \textbf{Parameter Analysis of the Most Probable Kinetic Scheme.}  
		In parallel with model comparison, we perform a detailed analysis of the parameters within the most probable kinetic scheme. In particular, the examination of kinetic coupling enables us to assess whether modulation of kinetic barriers occurs beyond simple regulation of equilibrium constants. Furthermore, evaluating the dependency of the current on the number of rotated subunits provides critical insight into how subunit rotation influences pore permeability.
	\end{enumerate}
	
	Collectively, these approaches provide a robust framework for deriving mechanistic interpretations from macroscopic current recordings and for elucidating the kinetics of P2X2 receptor activation.
	

\subsection{Likelihood Function Approximations of Macroscopic Currents}

\subsubsection{MacroIR and Related Algorithms}
In this section, we introduce \textbf{MacroIR}, a novel algorithm for determining the likelihood of time-averaged macroscopic currents. MacroIR adapts the equations of our previously published macroscopic recursive algorithm (\textbf{MacroR})~\cite{Moffatt} for calculating the posterior likelihood. However, instead of operating on the instantaneous state, MacroIR operates on a \emph{boundary state}—the composite state defined by the channel ensemble at both the beginning and the end of the time interval over which the current is averaged.

The rationale for focusing on the boundary state is rooted in the Markov property, which asserts that future dynamics depend solely on the present state rather than on the complete history. In our framework, the state at the beginning of an interval is sufficient to predict the dynamics within that interval, and the state at the end captures all the information necessary for forecasting the subsequent interval. To model the dynamics during the interval, we compute the mean of the averaged current conditioned on these boundary states and evaluate the variance of this mean. This strategy directs our analysis to the overall signal properties, eschewing the transient fluctuations that occur within the interval.

Previously, we introduced several algorithms for estimating the likelihood of instantaneous current measurements, including \textbf{MacroR} (macroscopic recursive), \textbf{MicroR} (microscopic recursive), and \textbf{MacroNR} (macroscopic non-recursive)~\cite{Moffatt}. MacroR iteratively updates the ensemble state of channels, akin to a Kalman filter, while MicroR applies Bayes’ rule to ensembles of channels with identical kinetic properties to compute the exact likelihood. Additionally, we introduce \textbf{SingleR}, a recursive method applied to individual (single-channel) states that redefines a well-established concept~\cite{Qin}.

To extend our framework to time-interval data using the concept of boundary states, we define two complementary microscopic methods. \textbf{MicroIR} forms the foundation of MacroIR by applying Bayesian inference to time-averaged measurements. \textbf{SingleIR} offers an alternative approach designed to address the missing event problem in single-molecule traces (though detailed demonstrations are not provided here). Collectively, these algorithms establish a comprehensive framework for analyzing current measurements across both microscopic and macroscopic scales.

\subsubsection{Notation}
We denote specific channel states using subscripts (e.g., \( i, j \)), while superscripts (e.g., ``prior", ``post", ``obs") indicate the type of probability distribution. Transitions from an initial state to an end state over an interval are denoted using arrows (e.g., $i \rightarrow j$). For instance, \( p_i^{\text{prior}} \) represents the prior probability of being in state \( i \), and $\Pi_{i \rightarrow j}^{\text{prior}}$ represents the prior probability of starting at state $i$ and ending at state $j$.

\subsubsection{Markov Model of Single-Channel Behavior}
\label{sec:SingleChannel}
Ion channel gating is inherently stochastic. We model channel kinetics as a Markov process with a finite set of states \(k\). The state occupancy is described by the probability vector \(\boldsymbol{p}(t)\), which evolves according to:
\begin{equation}
	\frac{d\boldsymbol{p}}{dt}(t) = \boldsymbol{p}(t)\cdot \boldsymbol{Q},
	\label{eq:master_equation}
\end{equation}
where the off-diagonal elements of \(\boldsymbol{Q}\) are the transition rates \(k_{ij}\) and the diagonal entries satisfy \( k_{ii} = -\sum_{j\neq i} k_{ij} \). The solution is given by:
\begin{equation}
	\boldsymbol{p}(t) = \boldsymbol{p}(0) \cdot \exp(\boldsymbol{Q}t) = \boldsymbol{p}(0) \cdot \mathbf{P}(t).
	\label{eq:master_solution}
\end{equation}
Since only the current generated by each state \(\gamma_i\) is measurable, the predicted instantaneous current is:
\begin{equation}
	y^{\text{pred}}(t) = \boldsymbol{p}(t)\cdot \boldsymbol{\gamma}.
	\label{eq:current_pred}
\end{equation}

\subsubsection{MicroR: Microscopic Bayesian Inference for Ensemble State Counts}
Experimental measurements often involve ensembles of \(N_{\text{ch}}\) identical channels. A full microscopic treatment requires tracking the probability distribution over the specific number of channels occupying each state, known as the ensemble's state configuration \( \boldsymbol{N} = (N_1, N_2, \dots, N_k) \), subject to \( \sum_i N_i = N_{\text{ch}} \).

We previously introduced the \textit{Microscopic Recursive} (MicroR) algorithm~\cite{Moffatt} to perform exact Bayesian inference on this ensemble. Given a prior probability distribution over configurations, \( P^{\text{prior}}(\boldsymbol{N}) \), and an observed instantaneous current \(y^{\text{obs}}\), MicroR calculates the likelihood by summing over all possible configurations:
\begin{equation}
	\mathcal{L} = P(y^{\text{obs}}) = \sum_{\boldsymbol{N}} P^{\text{prior}}(\boldsymbol{N}) P(y^{\text{obs}} | \boldsymbol{N}).
	\label{eq:microR_likelihood}
\end{equation}
Here, \( P(y^{\text{obs}} | \boldsymbol{N}) \) is typically modeled as \( \mathcal{N}(y^{\text{obs}}; \sum_i N_i \gamma_i, \epsilon^2) \). Bayes' theorem is then used to update the probability of each configuration:
\begin{equation}
	P^{\text{post}}(\boldsymbol{N}) = \frac{ P^{\text{prior}}(\boldsymbol{N}) P(y^{\text{obs}} | \boldsymbol{N})}{P(y^{\text{obs}})}.
	\label{eq:microR_posterior}
\end{equation}
While exact, the number of distinct configurations grows combinatorially with \( N_{\text{ch}} \), rendering MicroR computationally intractable for large macroscopic recordings. This necessitates the macroscopic approximations (MacroR and MacroIR) described subsequently.

\subsubsection{MicroIR: Microscopic Bayesian Inference for Time-Averaged Measurements}
Instantaneous observations are physically impossible; experiments measure time-averaged currents \(\overline{y}^{\text{obs}}\) over an interval \(t\). MicroIR extends Bayesian analysis to these averaged measurements. The time average is simply the integration of the current between the start ($t_n$) and end ($t_{n+1}$) of an interval.

MacroIR calculates the probability of a sequence of time averages given a model $M$ with parameters $\beta$. Due to the Markov property, to handle the interval $t_{n+1} \rightarrow t_{n+2}$, we require only the state at the start of that interval ($s_{t_{n+1}}$). The likelihood of the sequence is a product of conditional probabilities, where the posterior of the state $s_{t_{n+1}}$ is obtained by marginalizing over the state at the start of the previous interval $s_{t_n}$:
\begin{equation}
	P(y^{\text{obs}}_{t_{n}\rightarrow t_{n+1}} | s_{t_n}) = \sum_{j} P(y^{\text{obs}}_{t_{n}\rightarrow t_{n+1}} | s_{t_n}=i \rightarrow s_{t_{n+1}}=j) \cdot P(s_{t_{n+1}}=j | s_{t_n}=i).
\end{equation}
MicroIR relies on the joint probability of starting in state \(i\) and ending in state \(j\), defined as:
\begin{equation}
	\Pi_{i\rightarrow j}^{\text{prior}} = p_i^{\text{prior}}\,P_{i\rightarrow j}(t).
	\label{eq:joint_state_prob}
\end{equation}
The likelihood of observing $\overline{y}^{\text{obs}}$ is given by:
\begin{equation}
	\mathcal{L} = \sum_{i,j} \Pi_{i\rightarrow j}^{\text{prior}}\,\mathcal{N}\!\Biggl(
	\overline{y}^{\text{obs}}; \overline{\Gamma}_{i\rightarrow j}, \,
	\epsilon ^2_{0 \rightarrow t}  + \sigma^2_{\overline{\Gamma}_{i\rightarrow j}}
	\Biggr),
	\label{eq:integrated_likelihood}
\end{equation}
where $\overline{\Gamma}_{i\rightarrow j}$ and $\sigma^2_{\overline{\Gamma}_{i\rightarrow j}}$ are the mean and variance of the current for the transition $i \rightarrow j$. The noise term is defined as $\epsilon ^2_{0 \rightarrow t} = \frac{\epsilon^2}{t} + \nu ^2$, combining white noise $\epsilon$ and pink noise $\nu$.

Using Bayes' rule, the posterior joint probability is calculated, and the prior for the next interval is obtained by marginalization:
\begin{equation}
	p_j^{\text{prior}}(t_{n+1}) = \sum_i \Pi_{i\rightarrow j}^{\text{post}}.
	\label{eq:next_prior_microIR}
\end{equation}

\subsubsection{Calculation of Conditional Moments of Current}
\label{sec:moments}

To implement MicroIR (and subsequently MacroIR), we require the moments of the current conditional on the boundary states.

\paragraph{Mean Current:}
The time-averaged current for transitions from state \(i\) to state \(j\) over interval \(t\) is:
\begin{equation}
	\overline{\Gamma}_{i \rightarrow j} = \frac{1}{t\,P_{i\rightarrow j}} \int_0^t \sum_k P_{i\rightarrow k}(\tau)\,\gamma_k\,P_{k\rightarrow j}(t-\tau)\,d\tau.
	\label{eq:gamma_ij_integral}
\end{equation}
Using spectral decomposition of \(\mathbf{Q}\) (eigenvectors $\mathbf{V}$, eigenvalues $\lambda$), this yields a closed-form solution:
\begin{equation}
	\overline{\Gamma}_{i \rightarrow j} = \frac{1}{P_{i\rightarrow j}} \sum_{k, n_1, n_2} V_{i n_1}\,V^{-1}_{n_1 k}\,\gamma_k\,V_{k n_2}\,V^{-1}_{n_2 j}\,E_2(\lambda_{n_1}t,\lambda_{n_2}t),
	\label{eq:gamma_ij_formula}
\end{equation}
where
\begin{equation}
	E_2(x,y)= \begin{cases} \frac{e^x-e^y}{x-y}, & x\neq y, \\ e^x, & x=y. \end{cases}
\end{equation}

\paragraph{Variance of Current:}
The variance is $\mathrm{Var}(\overline{\Gamma}_{i\rightarrow j}) = E\bigl[\overline{\Gamma}_{i\rightarrow j}^2\bigr] - \Bigl(E[\overline{\Gamma}_{i\rightarrow j}]\Bigr)^2$. The expected square is:
\begin{equation}
	E\bigl[\overline{\Gamma}_{i\rightarrow j}^2\bigr] = \frac{1}{P_{i\rightarrow j}} \sum_{k_1,k_2,n_1,n_2,n_3} \left[ V_{i n_1}V^{-1}_{n_1 k_1}\gamma_{k_1}V_{k_1 n_2}V^{-1}_{n_2 k_2}\gamma_{k_2}V_{k_2 n_3}V^{-1}_{n_3 j} \right] E_3(\lambda_{n_1}t,\lambda_{n_2}t,\lambda_{n_3}t).
	\label{eq:expected_square_closed}
\end{equation}
The auxiliary function \(E_3\) is defined as:
\begin{equation}
	E_3(x,y,z)= 
	\begin{cases}
		\frac{e^x}{(x-y)(x-z)} + \frac{e^y}{(y-z)(y-x)} + \frac{e^z}{(z-x)(z-y)}, & x\neq y \neq z \neq x,\\[1mm]
		E_{12}(x,y),   & x\neq y, y=z,\\[1mm]
		E_{12}(y,z),   & y\neq z, z=x,\\[1mm]
		E_{12}(z,x),   & z\neq x, x=y,\\[1mm]
		\frac{1}{2}e^x, & x=y=z,
	\end{cases}
	\label{eq:E3}
\end{equation}
where
\begin{equation}
	E_{12}(x,y)= \frac{e^x}{(x-y)^2} + \frac{e^y}{y-x}\Bigl(\frac{y-x-1}{y-x}\Bigr).
\end{equation}

\subsubsection{MacroR: Macroscopic Bayesian Framework for Instantaneous Measurements}
Given the computational cost of MicroR, we developed \textbf{MacroR}~\cite{Moffatt} for instantaneous measurements. MacroR approximates the ensemble state distribution with a multivariate normal distribution defined by the mean state probability vector \( \boldsymbol{\mu} \) and the covariance matrix \( \boldsymbol{\Sigma} \).

The likelihood of observing an instantaneous current \( y^{\text{obs}} \) is:
\begin{equation}
	\mathcal{L} = \mathcal{N}\Bigl( y^{\text{obs}} ; y^{\text{pred}},\, \sigma^2_{y^{\text{pred}}} \Bigr),
	\label{eq:macro_likelihood}
\end{equation}
where the predicted current and its variance are:
\begin{align}
	y^{\text{pred}} &= N_{\text{ch}}\, \boldsymbol{\mu}^{\text{prior}} \cdot \boldsymbol{\gamma}, \\
	\sigma^2_{y^{\text{pred}}} &= \epsilon^2 + \nu^2 + N_{\text{ch}}\, (\boldsymbol{\gamma}^{\mathrm{T}}\, \boldsymbol{\Sigma}^{\text{prior}}\, \boldsymbol{\gamma}).
\end{align}
Here, $\nu^2$ represents residual model noise (1/f-like). The posterior update, mathematically equivalent to a Kalman filter~\cite{stepanyuk2011efficient}, is:
\begin{equation}
	\boldsymbol{\mu}^{\text{post}} = \boldsymbol{\mu}^{\text{prior}} + (y^{\text{obs}} - y^{\text{pred}})\, \frac{(\boldsymbol{\Sigma}^{\text{prior}} \boldsymbol{\gamma})^{\mathrm{T}}}{\sigma^2_{y^{\text{pred}}}},
\end{equation}
\begin{equation}
	\boldsymbol{\Sigma}^{\text{post}} = \boldsymbol{\Sigma}^{\text{prior}} - \frac{(\boldsymbol{\Sigma}^{\text{prior}}\, \boldsymbol{\gamma})\, (\boldsymbol{\Sigma}^{\text{prior}}\, \boldsymbol{\gamma})^{\mathrm{T}}}{\sigma^2_{y^{\text{pred}}}}.
\end{equation}
Propagation to the next time point ($t_{n+1}$) uses the transition matrix $\mathbf{P}$:
\begin{align}
	\boldsymbol{\mu}^{\text{prior}}(t_{n+1}) &= \boldsymbol{\mu}^{\text{post}}(t_n)\, \mathbf{P}, \\
	\boldsymbol{\Sigma}^{\text{prior}}(t_{n+1}) &= \mathbf{P}^{\mathrm{T}} \Bigl( \boldsymbol{\Sigma}^{\text{post}}(t_n) - \mathrm{diag}(\boldsymbol{\mu}^{\text{post}}(t_n)) \Bigr) \mathbf{P} + \mathrm{diag}(\boldsymbol{\mu}^{\text{prior}}(t_{n+1})).
\end{align}

\subsubsection{MacroIR: Macroscopic Bayesian Framework for Interval-Averaged Analysis}
\textbf{MacroIR} (Macroscopic Interval Recursive) extends the macroscopic Gaussian approximation to time-averaged measurements. It utilizes the \textbf{boundary state}—the joint probability of starting the interval \([0, t]\) in state \(i_0\) and ending in state \(i_t\).

While the boundary state implies a larger state space ($k \times k$), MacroIR avoids explicit computation of the full boundary covariance matrix. The predicted average current over the interval is:
\begin{equation}
	\overline{y}^{\text{pred}}_{0 \rightarrow t} = N_{\text{ch}} \cdot (\boldsymbol{\mu}^{\text{prior}}_{0} \cdot \overline{\boldsymbol{\gamma}}_{0}),
	\label{eq:macro_interval_predicted_y}
\end{equation}
where \( (\overline{\gamma}_{0})_i = \sum_j P_{i\rightarrow j}(t) \overline{\Gamma}_{i \rightarrow j} \) is the expected average current given the process starts in state \(i\).

The variance of this prediction includes instrumental noise, process noise, and initial state uncertainty:
\begin{equation}
	\sigma^2_{\overline{y}^{\text{pred}}_{0 \rightarrow t}} = \epsilon^2_{0 \rightarrow t} + N_{\text{ch}} \cdot (\boldsymbol{\mu}^{\text{prior}}_{0} \cdot \boldsymbol{\sigma}^2_{\overline{\gamma}_{0}}) + N_{\text{ch}} \cdot \widetilde{\mathbf{\gamma}^{\mathrm{T}} \mathbf{\Sigma}\mathbf{\gamma}},
	\label{eq:macro_interval_sigma_pred}
\end{equation}
where \( (\sigma^2_{\overline{\gamma}_{0}})_i = \sum_j P_{i\rightarrow j}(t) \mathrm{Var}(\overline{\Gamma}_{i\rightarrow j}) \). The term \( \widetilde{\mathbf{\gamma}^{\mathrm{T}} \mathbf{\Sigma}\mathbf{\gamma}} \) represents the variance contribution from the initial state distribution, computed efficiently as:
\begin{equation}
	\widetilde{\mathbf{\gamma}^{\mathrm{T}} \mathbf{\Sigma}\mathbf{\gamma}} = \overline{\boldsymbol{\gamma}}_{0}^{\mathrm{T}} \cdot \Bigl( \boldsymbol{\Sigma}^{\text{prior}}_{0} - \mathrm{diag}(\boldsymbol{\mu}^{\text{prior}}_0) \Bigr) \cdot \overline{\boldsymbol{\gamma}}_{0} 
	+ \boldsymbol{\mu}^{\text{prior}}_0 \cdot \Bigl( (\overline{\mathbf{\Gamma}}_{0 \rightarrow t} \circ \mathbf{P}(t)) \circ (\overline{\mathbf{\Gamma}}_{0 \rightarrow t} \circ \mathbf{P}(t)) \Bigr) \cdot \mathbf{1}. 
	\label{eq:simplified_boundary_state}
\end{equation}
Here, $\circ$ denotes the Hadamard (element-wise) product. The likelihood is then:
\begin{equation}
	\mathcal{L}_{0 \rightarrow t} = \mathcal{N}\Bigl( \overline{y}^{\text{obs}}_{0 \rightarrow t} - \overline{y}^{\text{pred}}_{0 \rightarrow t},\, \sigma^2_{\overline{y}^{\text{pred}}_{0 \rightarrow t}} \Bigr).
\end{equation}

To update the system for the next interval, we require the covariance between the interval average and the final state, marginalized over the initial states:
\begin{equation}
	(\widetilde{\boldsymbol{\gamma}^{\mathrm{T}} \boldsymbol{\Sigma}})_{i_t} = \sum_{i_0} \sum_{j_0, j_t} \overline{\Gamma}_{j_0 \rightarrow j_t} (\Sigma^{\text{prior}}_{0 \rightarrow t})_{(j_0 \rightarrow j_t)(i_0 \rightarrow i_t)}.
\end{equation}
This vector is used to update the prior covariance for the next interval directly:
\begin{multline}
	\boldsymbol{\Sigma}^{\text{prior}}(t) = \Biggl[ \mathbf{P}(t)^{\mathrm{T}} \Bigl( \boldsymbol{\Sigma}^{\text{prior}}_{0} - \mathrm{diag}(\boldsymbol{\mu}^{\text{prior}}_{0}) \Bigr) \mathbf{P}(t) + \mathrm{diag}(\boldsymbol{\mu}^{\text{prior}}(t)) \Biggr] \\ 
	- \frac{N_{\text{ch}}}{\sigma^2_{\overline{y}^{\text{pred}}_{0 \rightarrow t}}} \, (\widetilde{\boldsymbol{\gamma}^{\mathrm{T}} \boldsymbol{\Sigma}})^{\mathrm{T}} \cdot (\widetilde{\boldsymbol{\gamma}^{\mathrm{T}} \boldsymbol{\Sigma}}).
	\label{eq:prior_covariance_update}
\end{multline}
This formulation allows MacroIR to accurately process time-averaged data with computational efficiency comparable to instantaneous methods.

	\subsubsection{Likelihood Function and Model Parameters}
	The likelihood function defines the probability of observing the averaged recorded macroscopic current at each time interval, given a specific kinetic model and its associated parameters. The model parameters are categorized into distinct groups based on their functional role:
	
	\begin{itemize}
		\item \textbf{Kinetic Scheme Parameters:} Parameters that characterize the transitions between channel states, formulated either as state-based or allosteric schemes.
		\item \textbf{Channel Availability Model:} We modeled the evolution of the number of available channels over the course of the experiment by incorporating processes such as inactivation. In this model, two parameters are used: the total number of channels and the inactivation rate. Inactivation is represented as a distinct state that is irreversibly accessible from all other states at a uniform rate.
		\item \textbf{Conductance Model:} Parameters that define the single-channel conductance, including potential state-dependent variations. 
		\item \textbf{Measurement Noise Model:} A model that accounts for stochastic fluctuations in the recorded signal, capturing both white and pink noise components.
		\item \textbf{Baseline Current Estimation:} A parameter that quantifies the residual baseline current present in the recordings.
	\end{itemize}
	
	
	All model parameters have a direct physical or material interpretation, meaning they correspond to measurable or experimentally relevant quantities. This ensures that parameter estimates can be directly related to the underlying biological and experimental conditions.
	
	\subsubsection{Prior Distribution}
	For each model, we specify a prior distribution for every parameter. Rather than encoding detailed prior knowledge of purinergic receptor kinetics, our priors were designed to encompass a broad range of plausible kinetic models, constrained only by general considerations. In this way, we aimed to favor kinetic schemes that structurally reflect receptor dynamics.
	
	Since all parameters are inherently positive, we assigned each a log$_{10}$-normal prior with a common variance of 2 in log-space, ensuring that the bulk of the probability mass spans roughly two orders of magnitude around the median. The medians of these priors were set to reasonable values, avoiding undue bias toward any particular model configuration. In cases where symmetry is present, such as when two parameters are interchangeable (e.g., right versus left coupling or in schemes 3 and 4, which incorporate two equivalent open states), one branch was assigned a median 10 times greater than its counterpart.
	
	For the kinetic rate constants, we selected an intermediate value of 100 s$^{-1}$, while the association rate was set at 10 s$^{-1}$ mM$^{-1}$. The equilibrium allosteric couplings were assigned the following values: BF = 100, BG = 1, FG = 100, RB = 100, BR = 10, and RBR = 1. By contrast, all kinetic allosteric couplings were uniformly set to 1.
	
	Additional priors were specified as follows. The number of channels was fixed at 4800, the inactivation rate was set to $10^{-5}$ s$^{-1}$, and the measurement noise was defined with parameters of $10^{-3}$ pA$\cdot$s for white noise and 5 pA for pink noise. The unitary current was set to 1 pA (measured at $-60$ mV). Moreover, the leakage current was fixed at $10^{-6}$ pA, the allosteric factor (which increases the current due to subunit rotation) was set to 100, and the basal current was specified as 1 pA.   
	
	\subsubsection*{Likelihood}
	The likelihood is defined below (Equation~\eqref{eq:macroIR_likelihood}). Note that MacroIR is an approximation, valid only within certain regimes. In our implementation, if the approximation is found to be invalid over a given interval, we employ the MacroINR algorithm for that interval. This logic is implemented in the function \texttt{Macror} (line 3576 of \texttt{qmodel.h}; \url{https://github.com/lmoffatt/macro_dr_submission}).
	
	
	\subsection{Experimental Data: Outside-out Patch Recording}
	
	The experimental dataset analyzed in this study was previously published in \textit{The Journal of General Physiology} \cite{Moffatt_hume} and is reexamined here to develop and validate our approach. Recordings were obtained from macro outside-out patches containing hundreds to thousands of P2X2 channels. These patches were exposed to 0.2-ms ATP pulses at concentrations of 0.1, 0.2, 0.5, 1, 2, and 10 mM delivered every 2 minutes, interspersed with 10-ms pulses of 1 mM ATP. The dataset is ideally suited for kinetic analysis due to the precise alignment of stimulus timing with channel responses, which were recorded directly from the patch. For complete experimental details—including solution compositions, recording conditions, and data acquisition protocols—please refer to the original publication \cite{Moffatt_hume}.
	
	For kinetic analysis, we segmented the data as follows. The 1 s pre-pulse period was divided into 72 intervals of 13.7 ms and 10 intervals of 0.02 ms. During the ATP pulse, as well as 0.24 ms before,  the algorithm processed the raw data, whereas after the pulse the data were averaged over exponentially spaced intervals (approximately 10 points per decade). Additionally, for each 10-ms ATP pulse, one averaged sample  was retained from the second half of the pulse, corresponding to the equilibrium current. 
	
	\subsection{Kinetic Models Under Test}
	
	We evaluated 11 kinetic models, including seven previously described by Moffatt and Hume \cite{Moffatt_hume} and four new models introduced in this study. These models can be classified  into three distinct categories: \textbf{state models}, which describe receptor kinetics without explicitly incorporating allosteric interactions; \textbf{allosteric models}, which introduce modulations between specific transitions without an explicit structural correspondence; and \textbf{conformational models}, which explicitly map kinetic transitions onto physical structural rearrangements.
	
	\subsubsection{Previously Established Models}
	
	The seven previously established models represent progressively complex descriptions of receptor activation:
	
	\textbf{State Models:}  
	- \textit{Scheme I}: A Minimal linear scheme for a homotrimeric channel in which all binding sites must be occupied before channel opening can occur.  
	- \textit{Scheme II}: Extends Scheme I by introducing an intermediate "flipped" state (a hypothesized intermediate conformational state preceding opening) between the fully bound and open conformations, allowing for a sequential activation process.  
	- \textit{Scheme III}: Preferred scheme on a single channel study \cite{properties_single_channe}, this scheme introduces multiple open states converging into a final closed state, refining receptor deactivation pathways.  
	- \textit{Scheme IV}: Builds upon Scheme III by incorporating a flipped state.  
	
	\textbf{Allosteric Models:}  
	- \textit{Scheme V}: A \textbf{Monod-Wyman-Changeux (MWC)} model for a trimeric receptor, in which ligand binding promotes a concerted gating transition through an allosteric mechanism.  
	- \textit{Scheme VI}: Introduces an explicit \textbf{flipping transition} that can be allosterically modulated. In this scheme, binding increases the probability of flipping, which in turn increases the probability of opening.  
	- \textit{Scheme VII}: Extends Scheme VI by allowing flipping to occur independently at each binding site, creating a highly interconnected state network in which binding, flipping, and gating are coupled.  
	
	At the time these models were developed, it was not yet known that the ATP-binding site of P2X receptors was inter-subunit, as in Cys-loop receptors, rather than intra-subunit, as in glutamate receptors. Consequently, these models did not account for the complexity of ternary coupling. Later, the resolution of closed and open structures of zP2X4 \cite{cerrada_p2x, abierta_p2x} confirmed the inter-subunit nature of the ATP-binding site and suggested that agonist binding promotes subunit rotation, hinting at a potential structural mechanism for gating.
	
	A major limitation of the allosteric models (Schemes VI and VII) was that they not only postulated allosteric coupling between binding and flipping (and between flipping and gating)—which could be justified if flipping were the structural mechanism transmitting binding energy to the pore—but also required a direct binding–gating interaction to achieve a good fit. The latter was harder to interpret mechanistically. 
	
	\subsubsection{New Conformational Models}
	
	Inspired by structural insights from P2X4, we developed four new models that replace the hypothetical “flipping” state with a well‐defined \textbf{subunit partial rotation} and limit allosteric interactions to direct physical contacts—specifically, between ligand binding and subunit rotation. We considered three types of coupling based on the interacting elements:
	
	\begin{itemize}
		\item \textbf{RB (Rotation–Binding):} Coupling between the rotation of the subunit on the left of the ATP-binding site (subunit A in Figure~\ref{fig:distri}) and ATP binding.
		\item \textbf{BR (Binding–Rotation):} Coupling between the rotation of the subunit on the right of the ATP-binding site (subunit B in Figure~\ref{fig:distri}) and ATP binding.
		\item \textbf{RBR (Rotation–Binding–Rotation):} Ternary coupling that involves the rotation of both subunits forming the ATP-binding site in conjunction with ATP binding.
	\end{itemize}
	
	The four models are defined as follows:
	
	\begin{itemize}
		\item \textbf{Scheme VIII:} Incorporates only RBR coupling, where ATP binding induces a concerted rotation of both subunits that form the binding site.
		\item \textbf{Scheme IX:} Sets RBR to unity (i.e. no ternary coupling) and constrains RB equal to BR, thereby enforcing symmetric interactions between ligand binding and subunit rotations.
		\item \textbf{Scheme X:} Similar to Scheme IX, but allows RB and BR to vary independently, capturing potential asymmetries in the binding–rotation interactions.
		\item \textbf{Scheme XI:} Simultaneously incorporates all three coupling factors (RB, BR, and RBR), representing the most comprehensive model of receptor activation in this framework.
	\end{itemize}
	
	This conformational framework obviates the need for poorly defined allosteric interactions and provides a direct mapping between kinetic parameters and the structural changes observed in crystallographic studies.
	\subsubsection{Conductance Parameterization}
	
	For the conformational models (Schemes VIII, IX,X, XI), we hypothesized that conductance is a function of the number of rotated subunits. Instead of explicitly modeling the kinetics of a rotation–gating interaction (which would require a 48-state conformational model), we assumed that each rotated subunit stabilizes the open conformation, leading to an effective mean conductance that depends on the number of rotated subunits:
	
	\[
	i(n) = i_{\text{max}} \cdot \frac{E_n}{E_n + 1}, \quad E_n = E_0 \cdot F_g^n.
	\]
	
	Here, \( E_n \) represents the efficacy of the \( n \)-rotated state. This formulation significantly reduced the number of states while maintaining sufficient flexibility to capture conductance modulation. In addition to computational feasibility, this approach naturally accounts for the possibility that the pore lacks a strict blocking mechanism, and that its conductance is a nonlinear function of the number of rotated subunits.
	
	\subsubsection{Kinetic and Equilibrium Allosteric Factors}
	
	Previous studies have extended the allosteric framework—typically applied to equilibrium constants—to kinetic rates \cite{Moffatt_hume}. In this work, we refine this approach by explicitly defining kinetic allosteric factors.
	
	Consider two conformational transitions, \(g\) (gating) and \(r\) (rotation). In the absence of interaction, these transitions are characterized by four independent kinetic parameters: \(b_{\text{on}}, b_{\text{off}}, r_{\text{on}},\) and \(r_{\text{off}}\). When an allosteric interaction exists between binding and rotation (via either RB or BR coupling), the kinetic parameters in the activated state become
	\[
	b^*_{\text{on}},\quad b^*_{\text{off}},\quad r^*_{\text{on}},\quad r^*_{\text{off}}.
	\]
	
	Microscopic reversibility imposes the following constraint on these parameters:
	\[
	b_{\text{on}} \, r^*_{\text{on}} \, b^*_{\text{off}} \, r_{\text{off}} \;=\; r_{\text{on}} \, b^*_{\text{on}} \, r^*_{\text{off}} \, b_{\text{off}},
	\]
	which can be rearranged as:
	\[
	\frac{b_{\text{on}}/b_{\text{off}}}{b^*_{\text{on}}/b^*_{\text{off}}} \;=\; \frac{r_{\text{on}}/r_{\text{off}}}{r^*_{\text{on}}/r^*_{\text{off}}} \;=\; BR_{\text{eq}},
	\]
	for BR coupling (a similar relation holds for RB coupling). Thus, only two extra parameters are required to specify the coupled kinetic rates, defined in terms of the forward rates:
	\[
	BR_{b_{\text{on}}} \;=\; \frac{b^*_{\text{on}}}{b_{\text{on}}},
	\]
	\[
	BR_{r_{\text{on}}} \;=\; \frac{r^*_{\text{on}}}{r_{\text{on}}},
	\]
	with analogous definitions for RB coupling. For Bayesian model selection, we employed these forward-coupling parameters; however, for interpretative purposes, we also considered the corresponding reverse coupling factors:
	\[
	BR_{b_{\text{off}}} \;=\; \frac{b^*_{\text{off}}}{b_{\text{off}}},
	\]
	\[
	BR_{r_{\text{off}}} \;=\; \frac{r^*_{\text{off}}}{r_{\text{off}}}.
	\]
	\subsection{Likelihood function approximations of macroscopic currents}
	\subsubsection{MacroIR and Related Algorithms}
	In this section, we introduce \textbf{MacroIR}, a novel algorithm for determining the likelihood of time-averaged macroscopic currents. MacroIR adapts the equations of our previously published macroscopic recursive algorithm (\textbf{MacroR})~\cite{Moffatt} for calculating the posterior likelihood, repurposing them to operate on a \emph{boundary state}---the composite state defined by the channel ensemble at both the beginning and the end of the time interval over which the current is averaged---instead of on the instantaneous state.
	
	The rationale for focusing on the boundary state is rooted in the Markov property, which asserts that future dynamics depend solely on the present state rather than on the complete history. In our framework, the state at the beginning of an interval is sufficient to predict the dynamics within that interval, and the state at the end captures all the information necessary for forecasting the subsequent interval. To model the dynamics during the interval, we compute the mean of the averaged current conditioned on these boundary states and evaluate the variance of this mean, rather than the variance of the instantaneous current. This strategy directs our analysis to the overall signal properties, eschewing the transient fluctuations that occur within the interval.
	
	Previously, we introduced several algorithms for estimating the likelihood of instantaneous current measurements, including \textbf{MacroR} (macroscopic recursive), \textbf{MicroR} (microscopic recursive), and \textbf{MacroNR} (macroscopic non-recursive)~\cite{Moffatt}. MacroR iteratively updates the ensemble state of channels, akin to a Kalman filter, while MicroR applies Bayes’ rule to ensembles of channels with identical kinetic properties to compute the exact likelihood. In addition, we introduce \textbf{SingleR}, a recursive method applied to individual (single-channel) states that redefines a well-established concept~\cite{Qin}. Although SingleR and MicroR share the same underlying algebra, they differ in state definitions: SingleR targets individual channel events, whereas MicroR addresses ensemble behavior. MacroNR, originally presented by Milescu~\cite{Milescu} and later integrated into our framework~\cite{Moffatt}, provides a non-iterative alternative by estimating the likelihood without updating the channel probability distribution after each measurement.
	
	To extend our framework to time-interval data using the concept of boundary states, we define two complementary microscopic methods. Specifically, \textbf{MicroIR} forms the foundation of MacroIR by applying Bayesian inference to time-averaged measurements, while \textbf{SingleIR} offers an alternative approach designed to address the missing event problem and is proposed for applications to single channels or molecular motors (although detailed demonstrations of this application are not provided here). Collectively, these algorithms establish a comprehensive framework for analyzing current measurements across both microscopic and macroscopic scales.
	\subsubsection{Notation}
	We denote specific channel states using subscripts (e.g., \( i, j \)), while superscripts (e.g., ``prior", ``post", ``obs") are used to indicate the type of probability distribution (e.g., prior, posterior, or observed) and we denote the intervals and the transitions from the initial state to the end state occur using arrows (e.g., $\rightarrow$) . For instance, \( p_i^{\text{prior}} \) represents the prior probability of being in state \( i \), and \( p_i^{\text{post}} \) represents the posterior probability after an observation; $p_{i \rightarrow j}^{\text{prior}}$ and $p_{i \rightarrow j}^{\text{post}}$, on the other hand represents the prior  and posterior probabilities of starting at the state $i$ and ending at the state $j$. 
	
	\subsubsection{Markov Model of Single-Channel Behavior}\label{SingleChannel Behavior}
	
	Ion channel gating is inherently stochastic. We model channel kinetics as a Markov process with a finite set of states \(k\). The state occupancy is described by the probability vector \(\boldsymbol{p}(t)\), which evolves according to
	\begin{equation}
		\frac{d\boldsymbol{p}}{dt}(t) = \boldsymbol{p}(t)\cdot \boldsymbol{Q},
		\label{eq:master_equation_short}
	\end{equation}
	where the off-diagonal elements of \(\boldsymbol{Q}\) are the transition rates \(k_{ij}\) and the diagonal entries satisfy
	\begin{equation}
		k_{ii} = -\sum_{j\neq i} k_{ij}.
		\label{eq:Q_diag_short}
	\end{equation}
	The solution is given by
	\begin{equation}
		\boldsymbol{p}(t) = \boldsymbol{p}(0) \cdot \exp(\boldsymbol{Q}t).
		\label{eq:master equation solution}
	\end{equation}
	Since only the current generated by each state \(\gamma_i\) is measurable, the predicted current is
	\begin{equation}
		y^{\text{pred}}(t) = \boldsymbol{p}(t)\cdot \boldsymbol{\gamma}.
		\label{eq:current_pred_short}
	\end{equation}
	This concise formulation links the hidden Markov dynamics to the observable current.
	
	\subsubsection{MicroR: Microscopic Bayesian Inference for Ensemble State Counts}
	% Replaces the previous subsection 4 which used potentially misleading single-channel equations under the MicroR name.
	% This version defines MicroR consistent with Moffatt (20XX) as operating on ensemble meta-states.
	
	Experimental measurements often involve no single but ensembles of \(N_{\text{ch}}\) identical channels. A full microscopic treatment requires tracking the probability distribution over the specific number of channels occupying each state, known as the ensemble's meta-state or state configuration \( \boldsymbol{N} = (N_1, N_2, \dots, N_k) \), subject to the constraint \( \sum_i N_i = N_{\text{ch}} \).
	
	We previously introduced the \textit{Microscopic Recursive} (MicroR) algorithm \cite{Moffatt} to perform exact Bayesian inference on this ensemble state distribution using instantaneous current measurements. Given a prior probability distribution over configurations, \( P^{\text{prior}}(\boldsymbol{N}) \), and an observed instantaneous current \(y^{\text{obs}}\), MicroR calculates the likelihood by summing over all possible configurations:
	\begin{equation}
		\mathcal{L} = P(y^{\text{obs}}) = \sum_{\boldsymbol{N}} P^{\text{prior}}(\boldsymbol{N}) P(y^{\text{obs}} | \boldsymbol{N})
		\label{eq:microR_likelihood} % Changed label from potentially misleading single_channel_posterior
	\end{equation}
	Here, \( P(y^{\text{obs}} | \boldsymbol{N}) \) represents the probability of observing \(y^{\text{obs}}\) given the ensemble is in configuration \( \boldsymbol{N} \). Assuming Gaussian instrumental noise with variance \(\epsilon^2\), this is typically \( \mathcal{N}(y^{\text{obs}}; \sum_i N_i \gamma_i, \epsilon^2) \), where \( \sum_i N_i \gamma_i \) is the total current generated by configuration \( \boldsymbol{N} \). Bayes' theorem is then used to update the probability of each configuration:
	\begin{equation}
		P^{\text{post}}(\boldsymbol{N}) = \frac{ P^{\text{prior}}(\boldsymbol{N}) P(y^{\text{obs}} | \boldsymbol{N})}{P(y^{\text{obs}})}.
		\label{eq:microR_posterior} % Changed label
	\end{equation}
	Between measurements, the probability distribution \( P(\boldsymbol{N}, t) \) evolves according to the chemical master equation governing the ensemble dynamics.
	
	While MicroR provides an exact solution, the number of distinct configurations \( \boldsymbol{N} \) grows very rapidly with the number of channels \( N_{\text{ch}} \) and the number of states \(k\). Consequently, MicroR becomes computationally intractable for typical macroscopic recordings. This computational barrier necessitates the use of approximations, such as the macroscopic moment-based approaches (MacroR and MacroIR) described subsequently, which avoid explicitly tracking the probability of every configuration.
	
	% Note: This replaces the previous content under the MicroR heading. 
	% The original equations (5) and (6) seemed to describe single-channel updates and are removed. 
	% Original equations (7) and (8) described propagation and cumulative likelihood; 
	% the concepts still apply to P(N) but the detailed equations are omitted here for brevity.
	% Remember to replace \cite{Moffatt} with your actual bibtex key and ensure the citation appears correctly.
	\subsubsection{MicroIR: Microscopic Bayesian Inference for Time-Averaged Measurements}
	Instantaneous observations of ion channel currents are not technically possible, all experiments measure in fact time-averaged currents \(\overline{y}^{\text{obs}}\) over some interval \(t\). MicroIR extends classical Bayesian analysis of ensemble of channels to time-averaged measurements. El promedio temporal no es otra cosa que la integración entre dos tiempos, el inicio del intervalo y su finalización. MacroIR calcula la probabilidad de una serie de promedios temporales dado un modelo M con parámetros $\beta$ a partir de una serie de productos de probabilidades condicionales:
	$$
	P(y^{\text{obs}}_{t_{i}\rightarrow t_{i+1}},y^{\text{obs}}_{t_{i+1}\rightarrow t_{i+2}},\dots | M(\beta))=
	P(y^{\text{obs}}_{t_{i}\rightarrow t_{i+1}} | M(\beta)) \cdot P(y^{\text{obs}}_{t_{i+1}\rightarrow t_{i+2}}|y^{\text{obs}}_{t_{i}\rightarrow t_{i+1}}, M(\beta))\cdot \dots
	$$
	Ahora, como consecuencia de la propiedad de Markov, todo lo que necesitamos para el período $t_{n+1}\rightarrow t_{n+2}$ es el estado al inicio de este período: 
	$$
	P(y^{\text{obs}}_{t_{n+1}\rightarrow t_{n+2}}|y^{\text{obs}}_{t_{n}\rightarrow t_{n+1}}, M(\beta))= \sum_{i} P(y^{\text{obs}} _{t_{n+1}\rightarrow t_{n+2}}~|~s_{t_{n+1}}=i) \cdot P(s_{t_{n+1}}=i~|~y^{\text{obs}}_{t_{n}\rightarrow t_{n+1}}, M(\beta))
	$$
	Entonces tenemos que obtener el posterior de $s_{t_{n+1}}$ dada el promedio temporal $y^{\text{obs}}_{t_{n}\rightarrow t_{n+1}}$ marginalizando para el estado al inicio del intervalo $s_{t_n}$. 
	
	De modo que lo unico importante es el conocimiento del estado del sistema en las boundaries de los intervalos, lo que ocurre durante los intervalos es promediado. 
	
	lo cual nos indica que como 
	El meollo de la cuestión está en determinar la funcion recursiva 
	La propiedad de Markov indica que para 
	
	Upon some reflection on the Markov property, it is clear that o 
	
	
	Although the channel state fluctuates continuously during \(t\), Markov property indicates that all the information we need is the probability state at the beginning of the interval, the probability of the state trajectory can be determined using Equation~\ref{eq:master equation solution}
	at the  order to calculate the probability distribution of the averaged current over the successive intervals $t_{i}\rightarrow t_{i+1}$ and $t_{i+1}\rightarrow t_{i+2}$, we need not only the prior probability of the initial state $i$ but also the end state $i+1$ to calculate the for the first interval and the final state ofand for calculating the probability of the next interval  the end state $i_2$. 
	MacroIR uses the Normal distribution as an approximation to the probability distribution of the average current. Therefore, it has to calculate the mean current during the interval  during 
	MicroIR relies solely on the initial state probability \(\boldsymbol{p}^{\text{prior}}\) and its evolution via \(\mathbf{P}(t)\) (Eq.~\ref{eq:master equation solution}).
	
	The joint probability of starting in state \(i\) and ending in state \(j\) is defined as
	\begin{equation}
		\Pi_{i\rightarrow j}^{\text{prior}} = p_i^{\text{prior}}\,P_{i\rightarrow j},
		\label{eq:joint_state_probability_short}
	\end{equation}
	or equivalently, in matrix form, \(\boldsymbol{\Pi}^{\text{prior}} = \mathrm{diag}(\boldsymbol{p}^{\text{prior}})\,\mathbf{P}\).
	
	The likelihood of observing $\overline{y}^{\text{obs}}$ is given by
	\begin{equation}
		\mathcal{L} = \sum_{i,j} \Pi_{i\rightarrow j}^{\text{prior}}\,\mathcal{N}\!\Biggl(
		\overline{y}^{\text{obs}}; \overline{\Gamma}_{i\rightarrow j}, \,
		\epsilon ^2_{0 \rightarrow t}  + \sigma^2_{\overline{\Gamma}_{i\rightarrow j}}
		\Biggr),
		\label{eq:integrated_likelihood_short}
	\end{equation}
	where $\Pi_{i\rightarrow j}^{\text{prior}}$ denotes the prior probability of a transition from state $i$ to state $j$. The quantities $\overline{\Gamma}_{i\rightarrow j}$ and $\sigma^2_{\overline{\Gamma}_{i\rightarrow j}}$ represent the mean and variance, respectively, of the current associated with these transitions. The term $\epsilon ^2_{0 \rightarrow t}$ captures the variability not attributed to the model, which is decomposed into a white noise component $\epsilon^2$ and a pink noise component $\nu^2$ 
	\begin{equation}
		\epsilon ^2_{0 \rightarrow t} = \frac{\epsilon^2}{t} + \nu ^2
		\label{eq:interval_noise}
	\end{equation}
	
	
	Using Bayes' rule, the posterior joint probability becomes
	\begin{equation}
		\Pi_{i\rightarrow j}^{\text{post}} = \frac{\Pi_{i\rightarrow j}^{\text{prior}}\,\mathcal{N}\!\Bigl(\overline{y}^{\text{obs}}; \overline{\Gamma}_{i\rightarrow j},\, \epsilon ^2_{0 \rightarrow t} + \sigma^2_{\overline{\Gamma}_{i\rightarrow j}}\Bigr)}{P(\overline{y}^{\text{obs}})}.
		\label{eq:integrated_posterior_short}
	\end{equation}
	Marginalization over the initial state yields the prior for the next interval:
	\begin{equation}
		p_j^{\text{prior}}(t_{n+1}) = \sum_i \Pi_{i\rightarrow j}^{\text{post}}.
		\label{eq:next_prior_short}
	\end{equation}
	Cumulative log-likelihood is then updated additively over successive intervals. Although the exact distribution of the mean current is not strictly normal, we approximate it as such, assuming instrumental noise dominates over gating noise.
	
	This approach, while not applied in the present work, provides a conceptual basis for MacroIR and can be generalized to other single-molecule systems, such as molecular motors.
	
	
	\subsubsection{Calculation of the Average Current Conditional on Starting and Ending States}
	
	The time-averaged current for transitions from state \(i\) to state \(j\) over an interval \(t\) is defined as
	\begin{equation}
		\overline{\Gamma}_{i \rightarrow j} = \frac{1}{t\,P_{i\rightarrow j}} \int_0^t \sum_k P_{i\rightarrow k}(\tau)\,\gamma_k\,P_{k\rightarrow j}(t-\tau)\,d\tau,
		\label{eq:gamma_ij_integral_short}
	\end{equation}
	where \(P_{i\rightarrow j}\) is the overall transition probability and \(\gamma_k\) is the current associated with state \(k\).
	
	Employing the spectral decomposition of the rate matrix \(\mathbf{Q}\), this expression can be recast in closed form:
	\begin{equation}
		\overline{\Gamma}_{i \rightarrow j} = \frac{1}{P_{i\rightarrow j}} \sum_{k, n_1, n_2} V_{i n_1}\,V^{-1}_{n_1 k}\,\gamma_k\,V_{k n_2}\,V^{-1}_{n_2 j}\,E_2(\lambda_{n_1}t,\lambda_{n_2}t),
		\label{eq:gamma_ij_formula_short}
	\end{equation}
	where \(\mathbf{V}\) and \(\mathbf{V}^{-1}\) are the eigenvector matrix of \(\mathbf{Q}\) and its inverse, \(\lambda_{n}\) are the corresponding eigenvalues, and
	\begin{equation}
		E_2(x,y)=
		\begin{cases}
			\frac{e^x-e^y}{x-y}, & x\neq y, \\
			e^x, & x=y.
		\end{cases}
		\label{eq:E2_short}
	\end{equation}
	
	
	\subsubsection{Variance of the Average Current Conditional on Starting and Ending States}
	
	The variance of the time-averaged current for transitions from state \(i\) to \(j\) is defined as
	\begin{equation}
		\mathrm{Var}(\overline{\Gamma}_{i\rightarrow j}) = E\bigl[\overline{\Gamma}_{i\rightarrow j}^2\bigr] - \Bigl(E[\overline{\Gamma}_{i\rightarrow j}]\Bigr)^2,
		\label{eq:variance_gamma_short}
	\end{equation}
	with \(E[\overline{\Gamma}_{i\rightarrow j}]\) given by Eq.~\ref{eq:gamma_ij_integral_short}.
	
	The expected square of the average current is initially written as
	\begin{equation}
		E\bigl[\overline{\Gamma}_{i\rightarrow j}^2\bigr] = \frac{1}{t^2\,P_{i\rightarrow j}} \int_0^t\int_0^{t-\tau_1} \sum_{k_1,k_2} P_{i\rightarrow k_1}(\tau_1)\,\gamma_{k_1}\,P_{k_1\rightarrow k_2}(\tau_2)\,\gamma_{k_2}\,P_{k_2\rightarrow j}(t-\tau_1-\tau_2)\,d\tau_1\,d\tau_2.
		\label{eq:expected_square_integral_short}
	\end{equation}
	
	Using the spectral decomposition of \(\mathbf{Q}\), a closed-form expression is obtained:
	\begin{equation}
		E\bigl[\overline{\Gamma}_{i\rightarrow j}^2\bigr] = \frac{1}{P_{i\rightarrow j}} \sum_{k_1,k_2,n_1,n_2,n_3} V_{i n_1}\,V^{-1}_{n_1 k_1}\,\gamma_{k_1}\,V_{k_1 n_2}\,V^{-1}_{n_2 k_2}\,\gamma_{k_2}\,V_{k_2 n_3}\,V^{-1}_{n_3 j}\,E_3(\lambda_{n_1}t,\lambda_{n_2}t,\lambda_{n_3}t),
		\label{eq:expected_square_closed_short}
	\end{equation}
	where the auxiliary function \(E_3\) is defined as
	\begin{equation}
		E_3(x,y,z)= 
		\begin{cases}
			E_{111}(x,y,z), & x\neq y,\; y\neq z,\; z\neq x,\\[1mm]
			E_{12}(x,y),   & x\neq y,\; x\neq z,\; y=z,\\[1mm]
			E_{12}(y,z),   & y\neq z,\; y\neq x,\; z=x,\\[1mm]
			E_{12}(z,x),   & z\neq x,\; z\neq y,\; x=y,\\[1mm]
			\frac{1}{2}e^x, & x=y=z.
		\end{cases}
		\label{eq:E3_short}
	\end{equation}
	Here,
	\begin{equation}
		E_{12}(x,y)= E_{1,11}(x,y,y) + \frac{e^y}{y-x}\Bigl(\frac{y-x-1}{y-x}\Bigr),
		\label{eq:E12_short}
	\end{equation}
	and
	\begin{equation}
		E_{1,11}(x,y,z)= \frac{e^x}{(x-y)(x-z)}.
		\label{eq:E1_11_short}
	\end{equation}
	\begin{equation}
		E_{111}(x,y,z)= \frac{e^x}{(x-y)\cdot (x-z)} +\frac{e^y}{(y-z)\cdot (y-x)} \frac{e^z}{(z-x)\cdot (z-y)} 
		\label{eq:E111_short}
	\end{equation}
	
	
	\subsubsection{From Single-Channel to Ensemble Channel Analysis}
	
	For an ensemble of ion channels, the state probability vector initially yields a multinomial distribution for the state counts. However, measuring the total current refines the posterior distribution and introduces inverse correlations among states generating the same current. A detailed but computationally intensive method—the \textit{Microscopic Recursive Algorithm} \cite{Moffatt}—constructs an ensemble state vector by applying Bayes’ rule to every possible state count. For larger systems, we adopt a more efficient \textit{Macroscopic Approach} that approximates the state distribution with a multivariate normal, using the covariance matrix to capture the observation-induced correlations.
	
	\subsubsection{MacroR: Macroscopic Bayesian Framework for Instantaneous Ensemble Measurements}
	% Revised based on user feedback regarding origin, framing, and noise terms.
	
	Given the computational cost of the exact MicroR algorithm (Section X.X.X), we previously developed the \textit{Macroscopic Recursive} (MacroR) approach as an efficient alternative for analyzing instantaneous measurements from channel ensembles \cite{Moffatt}. MacroR operates on macroscopic variables: the mean state probability vector \( \boldsymbol{\mu} = \mathbb{E}[\boldsymbol{p}] \) and the covariance matrix of state probabilities \( \boldsymbol{\Sigma} = \text{Cov}(\boldsymbol{p}, \boldsymbol{p}) \), which describe the distribution of channels across states under a macroscopic approximation. 
	% Note: Depending on whether Sigma represents covariance of probabilities p_i or counts N_i, equations might slightly differ. Assuming probabilities based on context.
	
	The core idea is to approximate the ensemble's state distribution with a multivariate normal distribution. This allows for a computationally tractable Bayesian update procedure. Within this framework, the likelihood of observing an instantaneous current \( y^{\text{obs}} \) is given by a Gaussian distribution:
	\begin{equation}
		\mathcal{L} = \mathcal{N}\Bigl( y^{\text{obs}} ; y^{\text{pred}},\, \sigma^2_{y^{\text{pred}}} \Bigr)
		\label{eq:macro_likelihood}
	\end{equation}
	where \( y^{\text{pred}} \) is the predicted current based on the prior mean state probabilities:
	\begin{equation}
		y^{\text{pred}} = N_{\text{ch}}\, (\boldsymbol{\mu}^{\text{prior}} \cdot \boldsymbol{\gamma}),
		\label{eq:macro_predicted_y} % Assuming row vector mu, column vector gamma
	\end{equation}
	and the variance of the prediction, \( \sigma^2_{y^{\text{pred}}} \), incorporates multiple sources of variability:
	\begin{equation}
		\sigma^2_{y^{\text{pred}}} = \epsilon^2 + \nu^2 + N_{\text{ch}}\, (\boldsymbol{\gamma}^{\mathrm{T}}\, \boldsymbol{\Sigma}^{\text{prior}}\, \boldsymbol{\gamma}).
		\label{eq:macro_sigma_pred} % Kept original form, implicitly using N_ch factor for total variance
	\end{equation}
	Here, \(\epsilon^2\) represents the variance of the instrumental white noise component. \(\nu^2\) represents a $\frac{1}{f}$ noise component (akin to 'pink noise' in its scaling with averaging time), intended to capture residual stochasticity not explicitly described by the kinetic model, such as model inaccuracies or contributions from unmodeled channel types. The final term, \( N_{\text{ch}}\, (\boldsymbol{\gamma}^{\mathrm{T}}\, \boldsymbol{\Sigma}^{\text{prior}}\, \boldsymbol{\gamma}) \), reflects the variance arising from the statistical distribution of the \(N_{\text{ch}}\) channels across different current-carrying states, as captured by the prior covariance matrix \( \boldsymbol{\Sigma}^{\text{prior}} \) within the multivariate normal approximation.
	
	MacroR uses a Bayesian approach to update the estimated mean and covariance based on the observation. This approach was subsequently recognized as being mathematically equivalent to a Kalman filter formulation applied to this system \cite{stepanyuk2011efficient}. The posterior mean and covariance are updated as follows:
	\begin{equation}
		\boldsymbol{\mu}^{\text{post}} = \boldsymbol{\mu}^{\text{prior}} + (y^{\text{obs}} - y^{\text{pred}})\, \frac{(\boldsymbol{\Sigma}^{\text{prior}} \boldsymbol{\gamma})^{\mathrm{T}}}{\sigma^2_{y^{\text{pred}}}},
		\label{eq:macro_mean_posterior} % Assuming mu is row vector
	\end{equation}
	\begin{equation}
		\boldsymbol{\Sigma}^{\text{post}} = \boldsymbol{\Sigma}^{\text{prior}} - \frac{(\boldsymbol{\Sigma}^{\text{prior}}\, \boldsymbol{\gamma})\, (\boldsymbol{\Sigma}^{\text{prior}}\, \boldsymbol{\gamma})^{\mathrm{T}}}{\sigma^2_{y^{\text{pred}}}}.
		\label{eq:macro_cov_posterior} % Note: Removed the explicit N/Nch factor here as Sigma is Cov(p,p)
	\end{equation}
	The prior estimates for the next time point \( t_{n+1} = t_n + \Delta t_n \) are obtained by propagating the posterior estimates through the Markov process dynamics using the transition matrix \( \mathbf{P}(\Delta t_n) = \exp(\boldsymbol{Q}\Delta t_n) \):
	\begin{equation}
		\boldsymbol{\mu}^{\text{prior}}(t_{n+1}) = \boldsymbol{\mu}^{\text{post}}(t_n)\, \mathbf{P}(\Delta t_n),
		\label{eq:macro_mean_next_prior}
	\end{equation}
	\begin{equation}
		\boldsymbol{\Sigma}^{\text{prior}}(t_{n+1}) = \mathbf{P}(\Delta t_n)^{\mathrm{T}} \Bigl( \boldsymbol{\Sigma}^{\text{post}}(t_n) - \mathrm{diag}(\boldsymbol{\mu}^{\text{post}}(t_n)) \Bigr) \mathbf{P}(\Delta t_n) + \mathrm{diag}\Bigl(\boldsymbol{\mu}^{\text{prior}}(t_{n+1})\Bigr).
		\label{eq:macro_mean_next_cov} % Using user's original equation (17) structure
	\end{equation}
	MacroR thus provides a recursive method to estimate the evolving state distribution of a channel ensemble from instantaneous measurements, forming a basis for the analysis of time-averaged data presented next with MacroIR.
	
	% Remember to replace \cite{Moffatt} and \cite{Stepanyuk} with your bibtex keys.
	% Ensure vector/matrix dimensions and multiplication order (row vs column vectors for mu) are consistent throughout the manuscript.
	% Clarified the noise terms epsilon^2 and nu^2 based on user input.
	% Acknowledged Moffatt (origin) and Stepanyuk (Kalman equivalence).
	% Used original equations provided by user, assuming Sigma = Cov(p,p) and consistency checks.
	\subsubsection{MacroR: Bayesian Framework for Instantaneous Channel Ensemble Measurements}
	
	In the \textit{MacroR algorithm}, we treat measurements as instantaneous, focusing on the ensemble's state via its mean probability vector and covariance matrix. This approach—equivalent to a Kalman filter\cite{stepanyuk2011efficient}—yields the likelihood
	\begin{equation}
		\mathcal{L} = \mathcal{N}\Bigl( y^{\text{obs}} - y^{\text{pred}},\, \sigma^2_{y^{\text{pred}}} \Bigr),
		\label{eq:macro_likelihood}
	\end{equation}
	with the predicted current
	\begin{equation}
		y^{\text{pred}} = N_{\text{ch}}\, \boldsymbol{\mu}^{\text{prior}}\, \boldsymbol{\gamma},
		\label{eq:macro_predicted_y}
	\end{equation}
	and variance
	\begin{equation}
		\sigma^2_{y^{\text{pred}}} = \epsilon^2 + \nu^2+ N_{\text{ch}}\, \boldsymbol{\gamma}^{\mathrm{T}}\, \boldsymbol{\Sigma}^{\text{prior}}\, \boldsymbol{\gamma}.
		\label{eq:macro_sigma_pred}
	\end{equation}
	The posterior mean and covariance are updated as
	\begin{equation}
		\boldsymbol{\mu}^{\text{post}} = \boldsymbol{\mu}^{\text{prior}} + \frac{y^{\text{obs}} - y^{\text{pred}}}{\sigma^2_{y^{\text{pred}}}}\, \boldsymbol{\gamma}^{\mathrm{T}}\, \boldsymbol{\Sigma}^{\text{prior}},
		\label{eq:macro_mean_posterior}
	\end{equation}
	\begin{equation}
		\boldsymbol{\Sigma}^{\text{post}} = \boldsymbol{\Sigma}^{\text{prior}} - \frac{N}{\sigma^2_{y^{\text{pred}}}}\, \boldsymbol{\Sigma}^{\text{prior}}\, \boldsymbol{\gamma}\, \boldsymbol{\gamma}^{\mathrm{T}}\, \boldsymbol{\Sigma}^{\text{prior}}.
		\label{eq:macro_cov_posterior}
	\end{equation}
	Propagation to the next time point is given by
	\begin{equation}
		\boldsymbol{\mu}^{\text{prior}}(t_{n+1}) = \boldsymbol{\mu}^{\text{post}}(t_n)\, \mathbf{P}(\Delta t_n),
		\label{eq:macro_mean_next_prior}
	\end{equation}
	\begin{equation}
		\boldsymbol{\Sigma}^{\text{prior}}(t_{n+1}) = \mathrm{diag}\Bigl(\boldsymbol{\mu}^{\text{prior}}(t_{n+1})\Bigr) + \mathbf{P}(\Delta t_n)^{\mathrm{T}} \Bigl( \boldsymbol{\Sigma}^{\text{post}}(t_n) - \boldsymbol{\mu}^{\text{post}}(t_n) \Bigr) \mathbf{P}(\Delta t_n).
		\label{eq:macro_mean_next_cov}
	\end{equation}
	\subsubsection{MacroIR: A Bayesian Framework for Interval-Averaged Analysis of Channel Ensembles}
	% Revised explanations for key terms based on user feedback (marginalization of boundary state quantities).
	
	We now introduce MacroIR (Macroscopic Interval Recursive), our Bayesian framework designed for analyzing sequences of time-averaged measurements from channel ensembles, such as those obtained in macropatch recordings over successive time intervals \([t_n, t_{n+1}]\). MacroIR extends the macroscopic approach by incorporating information about the system's dynamics *within* each averaging interval.
	
	It leverages the concept of the **boundary state**—the joint probability of starting the interval \([0, t]\) in state \(i_0\) and ending in state \(i_t\). The prior mean vector \(\boldsymbol{\mu}^{\text{prior}}_{0 \rightarrow t}\) and covariance matrix \(\boldsymbol{\Sigma}^{\text{prior}}_{0 \rightarrow t}\) characterizing this boundary state distribution are derived from the initial state distribution \((\boldsymbol{\mu}^{\text{prior}}_0, \boldsymbol{\Sigma}^{\text{prior}}_0)\) and the transition dynamics \(\mathbf{P}(t) = \exp(\boldsymbol{Q}t)\):
	\begin{equation}
		(\mu^{\text{prior}}_{0 \rightarrow t})_{(i_0 \rightarrow i_t)} = (\mu^{\text{prior}}_0)_{i_0} \, P_{i_0 \rightarrow i_t}(t),
		\label{eq:boundary_mean_prior}
	\end{equation}
	\begin{multline}
		(\Sigma^{\text{prior}}_{0 \rightarrow t})_{(i_0 \rightarrow i_t)(j_0 \rightarrow j_t)} = P_{i_0 \rightarrow i_t}(t)\Bigl((\Sigma^{\text{prior}}_0)_{i_0,j_0} - \delta_{i_0,j_0}\, (\mu^{\text{prior}}_0)_{i_0}\Bigr) P_{j_0 \rightarrow j_t}(t) \\
		+ \delta_{i_0,j_0}\,\delta_{i_t,j_t}\, (\mu^{\text{prior}}_0)_{i_0}\, P_{i_0 \rightarrow i_t}(t).
		\label{eq:boundary_covariance_prior}
	\end{multline}
	A key feature of MacroIR is that, due to the structure of the Bayesian updates for this system, the final likelihood calculation and the recursive propagation steps are formulated to **avoid the explicit computation and storage of the full posterior boundary state covariance matrix**, \(\boldsymbol{\Sigma}^{\text{post}}_{0 \rightarrow t}\). This avoidance, arising from algebraic simplification and marginalization during the update derivation, is crucial for computational efficiency compared to methods requiring explicit state-space expansion for boundary states.
	
	The influence of the boundary path (start and end states) is implicitly captured through quantities like the mean current \(\overline{\Gamma}_{i \rightarrow j}\) and its variance \(\sigma^2_{\overline{\Gamma}_{i \rightarrow j}}\) conditional on the start and end states (defined in Sections~\ref{sec:avg_current_calc} and \ref{sec:var_current_calc})\footnote{Make sure these section labels exist or adjust references.}, which are pre-computed from the model kinetics and used in the likelihood.
	
	The predicted average current over the interval \([0, t]\) depends on the initial state mean \(\boldsymbol{\mu}^{\text{prior}}_0\) and the expected currents marginalized over the end states:
	\begin{equation}
		\overline{y}^{\text{pred}}_{0 \rightarrow t} = N_{\text{ch}} \cdot (\boldsymbol{\mu}^{\text{prior}}_{0} \cdot \overline{\boldsymbol{\gamma}}_{0}),
		\label{eq:macro_interval_predicted_y}
	\end{equation}
	where \( (\overline{\gamma}_{0})_i = \sum_j P_{i\rightarrow j}(t) (\overline{\Gamma})_{i \rightarrow j} \) represents the expected average current over the interval given the process started in state \(i\). The variance of this prediction is:
	\begin{equation}
		\sigma^2_{\overline{y}^{\text{pred}}_{0 \rightarrow t}} = \epsilon^2_{0 \rightarrow t} + N_{\text{ch}} \cdot \widetilde{\mathbf{\gamma}^{\mathrm{T}} \mathbf{\Sigma}\mathbf{\gamma}} + N_{\text{ch}} \cdot (\boldsymbol{\mu}^{\text{prior}}_{0} \cdot \boldsymbol{\sigma}^2_{\overline{\gamma}_{0}}),
		\label{eq:macro_interval_sigma_pred}
	\end{equation}
	with interval noise \( \epsilon^2_{0 \rightarrow t} = \epsilon^2/t + \nu^2 \) (Eq.~\ref{eq:interval_noise}) and the expected variance conditional on the start state \( (\sigma^2_{\overline{\gamma}_{0}})_i = \sum_j P_{i\rightarrow j}(t) \mathrm{Var}(\overline{\Gamma}_{i\rightarrow j}) \). The term \( \widetilde{\mathbf{\gamma}^{\mathrm{T}} \mathbf{\Sigma}\mathbf{\gamma}} \) represents the contribution to the predicted current's variance arising specifically from the uncertainty in the initial state distribution (\( \boldsymbol{\Sigma}^{\text{prior}}_{0} \)). While derived from operations involving the prior boundary state statistics (mean current vector \(\boldsymbol{\gamma}_{0 \rightarrow t}\) and covariance \(\boldsymbol{\Sigma}^{\text{prior}}_{0 \rightarrow t}\)), it is computed efficiently using the relationship:
	\begin{multline}
		% Using the user's provided formula:
		\widetilde{\mathbf{\gamma}^{\mathrm{T}} \mathbf{\Sigma}\mathbf{\gamma}} = \overline{\boldsymbol{\gamma}}_{0}^{\mathrm{T}} \cdot \Bigl( \boldsymbol{\Sigma}^{\text{prior}}_{0} - \mathrm{diag}(\boldsymbol{\mu}^{\text{prior}}_0) \Bigr) \cdot \overline{\boldsymbol{\gamma}}_{0} \\ % Note: Original had gamma_0_t here, needs checking. Assuming gamma_0 based on symmetry.
		+ \boldsymbol{\mu}^{\text{prior}}_0 \cdot \Bigl( (\overline{\mathbf{\Gamma}}_{0 \rightarrow t} \circ \mathbf{P}(t)) \circ (\overline{\mathbf{\Gamma}}_{0 \rightarrow t} \circ \mathbf{P}(t)) \Bigr) \cdot \mathbf{1}. 
		\label{eq:simplified_boundary_state} % Keep original label if preferred
	\end{multline}
	This calculation (Eq.~\ref{eq:simplified_boundary_state}) bypasses the need to explicitly form \( \boldsymbol{\Sigma}^{\text{prior}}_{0 \rightarrow t} \), depending only on the initial state covariance \( \boldsymbol{\Sigma}^{\text{prior}}_{0} \) and pre-calculated mean current terms.
	
	The likelihood for observing the average current \( \overline{y}^{\text{obs}}_{0 \rightarrow t} \) over the interval is then:
	\begin{equation}
		\mathcal{L}_{0 \rightarrow t} = \mathcal{N}\Bigl( \overline{y}^{\text{obs}}_{0 \rightarrow t} ; \overline{y}^{\text{pred}}_{0 \rightarrow t},\, \sigma^2_{\overline{y}^{\text{pred}}_{0 \rightarrow t}} \Bigr).
		\label{eq:macroIR_likelihood}
	\end{equation}
	The Bayesian update uses the observation \( \overline{y}^{\text{obs}}_{0 \rightarrow t} \) to refine the state estimate. The prior distribution for the *next* interval (starting at time \(t\)) is obtained efficiently. The key quantity driving the update is the vector \( \widetilde{\boldsymbol{\gamma}^{\mathrm{T}} \boldsymbol{\Sigma}} \), indexed by the end state \( i_t \). This vector represents the covariance between the interval-averaged current and the probability of ending in state \(i_t\). It is obtained by mathematically performing the matrix multiplication of the mean current vector associated with boundary transitions\footnote{Let \(\boldsymbol{\gamma}_{0 \rightarrow t}\) denote the vector of mean currents \( \overline{\Gamma}_{j_0 \rightarrow j_t} \), indexed by boundary pairs \((j_0 \rightarrow j_t)\).} (\(\boldsymbol{\gamma}_{0 \rightarrow t}\)) with the prior boundary state covariance matrix (\(\boldsymbol{\Sigma}^{\text{prior}}_{0 \rightarrow t}\)), and then **marginalizing the result by summing over all possible initial states \( i_0 \)**:
	\begin{equation}
		(\widetilde{\boldsymbol{\gamma}^{\mathrm{T}} \boldsymbol{\Sigma}})_{i_t} = \sum_{i_0} \Big( \boldsymbol{\gamma}_{0 \rightarrow t} \cdot \boldsymbol{\Sigma}^{\text{prior}}_{0 \rightarrow t} \Big)_{(i_0 \rightarrow i_t)} = \sum_{i_0} \sum_{j_0, j_t} (\overline{\Gamma})_{j_0 \rightarrow j_t} (\Sigma^{\text{prior}}_{0 \rightarrow t})_{(j_0 \rightarrow j_t)(i_0 \rightarrow i_t)}. 
		\label{eq:interval_gamma_sigma_explained} % Using explained label
	\end{equation}
	This marginalized covariance vector \( \widetilde{\boldsymbol{\gamma}^{\mathrm{T}} \boldsymbol{\Sigma}} \) is then used to compute the prior covariance for the next interval, \( \boldsymbol{\Sigma}^{\text{prior}}(t) \), incorporating the information gain from \( \overline{y}^{\text{obs}}_{0 \rightarrow t} \):
	\begin{multline}
		\boldsymbol{\Sigma}^{\text{prior}}(t) = \Biggl[ \mathbf{P}(t)^{\mathrm{T}} \Bigl( \boldsymbol{\Sigma}^{\text{prior}}_{0} - \mathrm{diag}(\boldsymbol{\mu}^{\text{prior}}_{0}) \Bigr) \mathbf{P}(t) + \mathrm{diag}(\boldsymbol{\mu}^{\text{prior}}(t)) \Biggr]_{\text{Propagated Covariance}} \\ 
		- \frac{1}{\sigma^2_{\overline{y}^{\text{pred}}_{0 \rightarrow t}}} \, (\widetilde{\boldsymbol{\gamma}^{\mathrm{T}} \boldsymbol{\Sigma}})^{\mathrm{T}} \cdot (\widetilde{\boldsymbol{\gamma}^{\mathrm{T}} \boldsymbol{\Sigma}}). 
		\label{eq:prior_covariance_update_0_t} % Corrected N->Nch implicitly/by definition of Sigma. This is Eq 29.
	\end{multline}
	The prior mean for the next interval, \( \boldsymbol{\mu}^{\text{prior}}(t) \), is similarly updated (details omitted for brevity, but involves \( \widetilde{\boldsymbol{\gamma}^{\mathrm{T}} \boldsymbol{\Sigma}} \) and marginalization of the formal posterior boundary state mean given by Eq.~\ref{eq:macro_interval_posterior_mean}). 
	% Optional: Add back Eq 27 if needed, but focus is on the final recursive update structure for time t.
	% \begin{equation} \mathbf{\mu}^{\text{post}}_{0 \rightarrow t} = ... \label{eq:macro_interval_posterior_mean} \end{equation}
	
	In comparison to other interval-based analysis methods like \cite{Munch2022}, MacroIR potentially offers higher accuracy. By incorporating the interval average centered around each time point for the updates, the state estimate at time \(t\) implicitly benefits from information regarding transitions influencing both the interval ending at \(t\) and the interval starting at \(t\). This contrasts with methods applying only a single update based on the measurement within one interval \cite{Munch2022}. While MacroIR involves more computations per time step to achieve this, its computational complexity scaling per step is expected to be similar to approaches like \cite{Munch2022} that also avoid state-space expansion, because MacroIR efficiently bypasses the computation of the full posterior boundary state covariance.
	
	In summary, MacroIR provides a recursive Bayesian framework tailored for interval-averaged measurements. By exploiting the mathematical structure to avoid explicit calculation of posterior boundary state covariances—using efficient computations involving marginalization over initial states (Eqs.~\ref{eq:simplified_boundary_state}, \ref{eq:interval_gamma_sigma_explained}, \ref{eq:prior_covariance_update_0_t})—it computes likelihoods and propagates the state distribution, offering a potentially more accurate approach for analyzing macroscopic currents.
	
	% TODO: Ensure section labels for avg/var current are correct.
	% TODO: Double-check Eq 25 (\widetilde{gamma^T Sigma gamma}) term involving gamma_0 vs gamma_0_t.
	% TODO: Verify the footnote defining gamma_0->t is sufficient or if it needs formal definition earlier.
	% TODO: Replace \cite{Munch2022} with correct bibtex key.
	% TODO: Ensure definition/scaling of Sigma and Nch is handled consistently.\subsubsection{MacroIR: A Bayesian Framework for Interval-Averaged Analysis of Channel Ensembles}
	
	We introduce MacroIR (Macroscopic Interval Recursive), a Bayesian framework for analyzing time-averaged Markovian processes over successive intervals, as encountered in macropatch recordings. In contrast with previous time-averaging algorithms \cite{Munch2022}, 
	MacroIR leverages the \textbf{boundary state}, which represents the joint probability of starting in state \(i_0\) and ending in state \(i_t\) over an interval \([0,t]\).
	
	The boundary state prior mean is defined as:
	\begin{equation}
		(\mu^{\text{prior}}_{0 \rightarrow t})_{(i_0 \rightarrow i_t)} = (\mu^{\text{prior}})_{i_0} \, P_{i_0 \rightarrow i_t}(t),
		\label{eq:boundary_mean_prior}
	\end{equation}
	with the corresponding covariance:
	\begin{multline}
		(\Sigma^{\text{prior}}_{0 \rightarrow t})_{(i_0 \rightarrow i_t)(j_0 \rightarrow j_t)} = P_{i_0 \rightarrow i_t}\Bigl((\Sigma^{\text{prior}}_0)_{i_0,j_0} - \delta_{i_0,j_0}\, (\mu^{\text{prior}}_0)_{i_0}\Bigr) P_{j_0 \rightarrow j_t} \\
		+ \delta_{i_0,j_0}\,\delta_{i_t,j_t}\, (\mu^{\text{prior}}_0)_{i_0}\, P_{i_0 \rightarrow i_t}
		\label{eq:boundary_covariance_prior}
	\end{multline}
	
	The predicted average current over the interval is given by
	\begin{equation}
		\overline{y}^{\text{pred}}_{0 \rightarrow t} = N_{\text{ch}} \cdot \mathbf{\mu}^{\text{prior}}_{0} \cdot \overline{\gamma}_{0},
		\label{eq:macro_interval_predicted_y}
	\end{equation}
	where the marginalized current per state is
	\begin{equation}
		(\overline{\gamma}_{0})_i = \sum_j (\overline{\Gamma})_{i \rightarrow j}.
	\end{equation}
	
	The variance of the predicted current is:
	\begin{equation}
		\sigma^2_{\overline{y}^{\text{pred}}_{0 \rightarrow t}} = \epsilon^2_{0 \rightarrow t} + N_{\text{ch}} \cdot \widetilde{\mathbf{\gamma}^{\mathrm{T}} \mathbf{\Sigma}\mathbf{\gamma}} + N_{\text{ch}} \cdot \mathbf{\mu}^{\text{prior}}_{0} \cdot \sigma^2_{\overline{\gamma}_{0}},
		\label{eq:macro_interval_sigma_pred}
	\end{equation}
	with
	\begin{equation}
		(\sigma^2_{\overline{\gamma}_{0}})_i = \sum_j (\sigma^2\,\overline{\Gamma})_{i \rightarrow j}.
	\end{equation}
	
	The effective term \( \widetilde{\mathbf{\gamma}^{\mathrm{T}} \mathbf{\Sigma}\mathbf{\gamma}} \) is computed as:
	\begin{multline}
		\widetilde{\mathbf{\gamma}^{\mathrm{T}} \mathbf{\Sigma}\mathbf{\gamma}} = \mathbf{\gamma}^{\mathrm{T}}_{0 \rightarrow t} \cdot \mathbf{\Sigma}^{\text{prior}}_{0 \rightarrow t} \cdot \mathbf{\gamma}_{0 \rightarrow t} = \\
		\overline{\mathbf{\gamma}}_{0}^{\mathrm{T}} \cdot \Bigl( \mathbf{\Sigma}^{\text{prior}}_{0} - \mathrm{diag}(\mathbf{\mu}^{\text{prior}}_0) \Bigr) \cdot \overline{\mathbf{\gamma}}_{0 \rightarrow t} \\
		\quad + \mathbf{\mu}^{\text{prior}}_0 \cdot \Bigl( \overline{\mathbf{\Gamma}}_{0 \rightarrow t} \circ \overline{\mathbf{\Gamma}}_{0 \rightarrow t} \circ \mathbf{P} \Bigr) \cdot \mathbf{1},
		\label{eq:simplified_boundary_state}
	\end{multline}
	where \( \circ \) denotes the element-wise (Hadamard) product and \( \mathbf{1} \) is a vector of ones.
	
	The likelihood function for the time interval $0 \rightarrow t$ is therefore
	\begin{equation}
		\mathcal{L}_{0 \rightarrow t} = \mathcal{N}\Bigl( \overline{y}^{\text{obs}}_{0 \rightarrow t} - \overline{y}^{\text{pred}}_{0 \rightarrow t},\, \sigma^2_{\overline{y}^{\text{pred}}_{0 \rightarrow t}} \Bigr)
		\label{eq:macroIR_likelihood}
	\end{equation}
	
	Posterior updates are performed without explicitly computing the boundary state. The posterior mean is updated as:
	\begin{equation}
		\mathbf{\mu}^{\text{post}}_{0 \rightarrow t} = \mathbf{\mu}^{\text{prior}}_{0 \rightarrow t} + \frac{y^{\text{obs}}_{0 \rightarrow t} - y^{\text{pred}}_{0 \rightarrow t}}{\sigma^2_{\overline{y}^{\text{pred}}_{0 \rightarrow t}}} \cdot \widetilde{\boldsymbol{\gamma}^{\mathrm{T}} \boldsymbol{\Sigma}},
		\label{eq:macro_interval_posterior_mean}
	\end{equation}
	with
	\begin{equation}
		\widetilde{\boldsymbol{\gamma}^{\mathrm{T}} \boldsymbol{\Sigma}} = \mathbf{1}_0^{\mathrm{T}} \cdot \mathbf{\gamma}^{\mathrm{T}}_{0 \rightarrow t} \cdot \mathbf{\Sigma}^{\text{prior}}_{0 \rightarrow t}.
		\label{eq:interval_gamma_sigma}
	\end{equation}
	
	Finally, the prior covariance for the next interval is updated by marginalizing over the initial states:
	\begin{multline}
		\boldsymbol{\Sigma}^{\text{prior}}_{t} = \mathbf{P}^{\mathrm{T}} \Bigl( \boldsymbol{\Sigma}^{\text{prior}}_{0} - \mathrm{diag}(\boldsymbol{\mu}^{\text{prior}}_{0}) \Bigr) \mathbf{P} + \mathrm{diag}(\boldsymbol{\mu}^{\text{prior}}_{0} \cdot \mathbf{P}) \\
		- \frac{N}{\sigma^2_{\overline{y}^{\text{pred}}_{0 \rightarrow t}}} \, \widetilde{\boldsymbol{\gamma}^{\mathrm{T}} \boldsymbol{\Sigma}}^{\mathrm{T}} \cdot \widetilde{\boldsymbol{\gamma}^{\mathrm{T}} \boldsymbol{\Sigma}}.
		\label{eq:prior_covariance_update_0_t}
	\end{multline}
	
	In summary, MacroIR provides a recursive Bayesian framework for interval-averaged measurements, efficiently computing likelihoods and posterior distributions while circumventing the computational burden of explicit boundary state calculations.
	
	\subsection{Bayesian Model Comparison}  
	We compared 11 alternative kinetic schemes to evaluate distinct mechanistic hypotheses for P2X2 receptor activation. For each scheme, a fully specified Bayesian model was constructed with explicit likelihood functions and prior distributions as described above. Bayesian evidence was estimated using a modified parallel-tempering, affine-invariant\cite{goodman2010ensemble}Markov chain Monte Carlo (MCMC) algorithm  with dynamic temperature selection \cite{vousden2016dynamic}. Rather than enforcing a uniform exchange rate between neighboring temperature chains, our implementation aims to maintain a uniform product of $\Delta\log\mathcal{L}$ and $\Delta\beta$. This strategy equalizes exchange rates at high $\beta$, while reducing the frequency of extremely negative $\Delta\beta\cdot\log\mathcal{L}$ values at low $\beta$, thereby obviating the need to justify their exclusion. The algorithm efficiently sampled high-dimensional posterior distributions by dynamically adapting the temperature configuration during sampling, which minimized potential sampling biases and ensured a uniform error rate in log-evidence estimation across temperatures.
	
	\subsubsection{MCMC Convergence and Resolution}  
	MCMC convergence was assessed by monitoring the potential scale reduction factor (PSRF) for each parameter and derived quantity, ensuring that effective sample sizes (ESS) were sufficiently large. Detailed PSRF and ESS statistics, along with credibility intervals for log-evidence estimates, the number of iterations, computational time, and system specifications (number of cores and processor used), are provided in Supplementary Table~1. Supplementary Table~2 examines the convergence of parameter estimates for Scheme 10 using the MacroIR approach, while Supplementary Table~3 presents parameter estimates for \(\beta=0\), validating our implementation of the affine-invariant ensemble Monte Carlo Markov Chain by confirming the expected median and variance of the prior distributions.
	\subsubsection{Validation of MacroIR Using Simulated Data}  
	To assess the robustness of our Bayesian framework, we validated the MacroIR method using simulated recordings. Synthetic ATP-activated currents were generated based on the fitted kinetic parameters for Scheme X, ensuring that simulation conditions closely matched the experimental setup. We then applied MacroIR to this dataset to evaluate its ability to recover the original parameters. As shown in Supplementary Table~4, all but two parameters were successfully recovered within their respective credibility intervals. The two outliers were the estimated Number of Channels (overestimated) and the Single-Channel Current (underestimated).  
	
	These two parameters primarily influence stochastic fluctuations without affecting time kinetics. A larger estimated channel count reduces fluctuations through random averaging, leading to a compensatory underestimation of single-channel conductance to maintain the same overall current. Since MacroIR is an approximation, some degree of bias is expected. Here, the overestimation of the Number of Channels suggests a form of overfitting, where an underestimation of stochastic fluctuations leads to a compensatory overestimation of kinetic components. Nevertheless, with only two outliers among 16 parameters—each with credibility intervals covering 90\% of the total probability—these deviations fall within expected variability but warrant further investigation.
	\subsubsection{Posterior Predictive Checks}  
	To evaluate model fit, we performed posterior predictive checks. In the MacroIR approach, each measurement is modeled as a normal distribution with a predicted mean current and variance. These predicted distributions were compared to the experimental data by computing a log-likelihood measure; the agreement between the observed and expected log-likelihood values confirms that the model accurately reproduces the experimental data.
	
	\subsection{Conformational characterization of the closed form of P2X perturbed by the binding of an ATP in one binding site}
	
	\subsubsection{Models for Molecular Dynamics simulations}
	\label{sec:model_const}
	
	P2X receptors are homomeric or heteromeric assemblies of three monomers labeled P2X${_1}$ through P2X${_7}$. Up to date, Protein Data Bank includes experimental structures from distinct species of P2X${1}$ \cite{bennetts2024structural}, P2X${3}$ \cite{cristal_p2x3, wang2018,li2019,thach2025mechanistic,kim2024discovery}, P2X${4}$ \cite{abierta_p2x,cerrada_p2x,shen2023structural,shi2025human}, and P2X${7}$ \cite{mccarthy2019full,binding_zn_2016,karasawa2016,sheng2024structural,oken2024high,oken2024p2x7}. Most of the structures lack complete intracellular residues, which are crucial for conducting stable MD simulations and preventing channel drift \cite{piedominici_p2x_1}. Besides, when AlphaFold \cite{af3} is implemented to predict these regions, results display very low confidence.  
	
	
	P2X receptors vary in the extent of their permeation selectivity \cite{samways2014principles}. However their overall architecture —particularly the extracellular region— is conserved across subtypes and species. 
	Each channel resembles a chalice with each monomer having a dolphin-like shape. The extracellular domain is relatively large with binding sites located in the upper part, approximately 40 Å away from the membrane surface.
	This conserved structural design strongly implies a unified closed$\rightarrow$open pathway triggered by ATP.
	Only P2X${4}$ and P2X${7}$ were determined in both their closed and open states P2X${4}$ \cite{abierta_p2x,cerrada_p2x,mccarthy2019full,karasawa2016,sheng2024structural}. According to our experience in practicing reliable MD simulations \cite{piedominici_p2x_1,racigh2021, pierdominici2019charge,racigh2019positively}, and due to the common opening mechanism of the distinct subtypes, we examined the perturbations caused by ATP docking to a single binding site considering three distinct models based on zebrafish P2x${_4}$. Two of them implied the open (P2X$^{open}$) and closed (P2X$^{closed}$) form of the channel. These structures were used to define a vector that describes the changes in the extracellular vestibule between both conformations. This region
	presents the largest differences in the transition P2X$^{closed}$ $\xrightarrow{}$ P2X$^{open}$ \cite{abierta_p2x} and coupled to the ion conducting transmembrane domain  (See Fig. \ref{fig:distri}). 
	The models were built considering the crystal structures of the zebrafish P2X4 (Uniprot code: F8W463) available in the Protein Data Bank: codes 3I5D (closed) and 4DW1 (open). The following paragraphs explain the modifications made to the original PDB files to equalize the number of residues of both forms.
	In PDB code 3I5D, chain-A spans from Arg36 to Ile359, chain-B from Leu34 to Leu361, and chain-C from Leu34 to Trp358. To standardize the lengths, we added residues Leu34, Asn35, Val360, and Leu361 to chain-A and residues 359-361 to chain-C, using chain-B as a reference. PDB code 4DW1, by the other side, includes only one of the three monomers and it includes the sequence Arg36-Ile359. We supplemented this chain with residues Leu34, Asn35, 360, and 361 to match the total number of residues in the closed structure. The atomic coordinates for the other two chains were derived using symmetry information from the crystallographic data.
	Additionally, we incorporated residues Val31, Gly32, and Thr33 into each chain of both structures. The conformation of these inserted segments was predicted using AlphaFold3 \cite{af3}, which suggested that the new residues would continue the $\alpha$-helix motif of the previous residues.
	In this way, the final forms of P2X$^{closed}$ and P2X$^{open}$ contains residues from Val31 to Leu361 in each of their three chains. P2X$^{open}$, besides, contains the coordinates of an ATP molecule in each of its three binding site.
	
	
	The third model was constructed to investigate, through Molecular Dynamics (MD) simulations, the initial conformational changes caused by the interaction of an ATP molecule to the binding site conformed by chains A and B. Chain A refers to the one that contacts ATP via its head and left flipper subdomains (marked in blue in Fig. \ref{fig:distri}), while chain B is the one that does it through its dorsal fin (coloured in red in Fig. \ref{fig:distri}).
	To construct this model, an ATP molecule was docked into the binding site of P2X$^{closed}$ formed by chains A and B. 
	The ATP coordinates were determined by aligning the C$\alpha$ atom positions of P2X$^{open}$ to those of P2X$^{closed}$. This resulting model will be referred as P2X$^{closed}_{1-ATP}$. 
	
	\subsubsection{Vector definition for the close \texorpdfstring{$\longrightarrow$}{->} open transition of the extracellular vestibule}
	
	
	It has been described that the opening of the transmembrane channel involves an outward ‘flexing’ of the lower region of the extracellular vestibule (EV) \cite{abierta_p2x}. In order to measure how the binding of ATP affects the EV dynamical behaviour, we defined a vector for each of the two subunits that composed the binding site where the ATP was docked in P2X$^{closed}_{1-ATP}$. To construct these vectors, P2X$^{open}$ and P2X$^{closed}$ were first aligned and then the
	coordinates of C$\alpha$ atoms from residues Asp59-Ser65, Asn195-Lys205, Arg261-Pro275, and Phe327-Phe333 of P2X$^{closed}$ were subtracted from those corresponding to  P2X$^{open}$. 
	This was performed for both chains A and B. These transition vectors, P2X$^{closed}_A$ $\rightarrow$ P2X$^{open}_A$ and P2X$^{closed}_B$ $\rightarrow$ P2X$^{open}_B$, have their origin in the closed conformation and contain the information for the changes of the EV from P2X$^{closed}$ and P2X$^{open}$ for the selected chains. 
	
	\subsubsection{Molecular Dynamics simulations}
	\label{sec:met_md_simulations}
	
	Standard MD simulations were performed on P2X$^{closed}_{1-ATP}$ considering the following protocol.
	Initially, the system was fed into the Leap module of AMBER22 where it was first neutralized. Then Na$^+$ and Cl$^-$ ions were incorporated to produce a 0.15M concentration. Water molecules were added to form an octahedral cell, and special care was devoted to waters molecules located inside the protein in order to conserve the shape of the receptor during the simulations. 
	Otherwise the protein scaffold may shrink.
	The Amber19SB force field \cite{tian2019ff19sb} was employed for computing the potential energy of the protein and water molecules while the Lipid14 force field was chosen for POPC \cite{lipid14}. 
	The PMEMD.cuda modules of AMBER18 package was employed to run the simulations \cite{amber18}.
	
	This initial structure was first minimized at constant volume and then heated at NVT conditions from 0 K to 100 K during 500 ps. Next, we switched to NTP conditions, and the heating was continued until 303 K was reached. Harmonic restraints of 1.5 kcal/mol{\AA}$^2$ were applied on the C$_\alpha$ atoms of the protein in the two heating periods. 
	This was followed by four consecutive 10.0 ns simulations at 303 K in which the restraints were gradually diminished (0.5, 0.1, 0.05 and 0.01 kcal/mol {\AA}$^2$). A final equilibration stage of 200 ns was performed.
	The production phase consists of ten independent MD simulations started from the final structure of the equilibration stage. Each of these trajectories were run during 200 ns, conformations were sampled every xx ps, and its initial velocities were randomly chosen from a Maxwellian distribution at 303K.
	
	The SHAKE algorithm was implemented to constrain the bonds involving hydrogen atoms. Accordingly, the time step could be set to 2.0 fs.
	A cutoff radius of 10.0 {\AA} was imposed to the electrostatic interactions and the Particle Mesh Ewald method was employed to calculate the long-range Coulombic force \cite{ewald1,ewald2}.
	
	\textcolor{red}{check convergence}
	
	
	
	
	\subsection{Data Availability}  
	The output from the Monte Carlo Markov Chain sampling of the kinetic models used to generate all figures is available via Zenodo (DOI: \url{https://doi.org/10.5281/zenodo.15085038}). Processed experimental data can be accessed at \url{https://github.com/lmoffatt/macro_dr_submission/experiments}.
	
	\subsection{Code Availability}  
	Custom software implemented in C++20, including 11 kinetic schemes, MacroIR, MacroINR, and evidence estimation algorithms, as well as R Markdown scripts for processing MCMC outputs and reproducing figures, is available at \url{https://github.com/lmoffatt/macro_dr_submission}.
	
	
	
	\section*{Methods}
	
	\subsection*{Synthetic models and data}
	We simulate $k$-state CTMCs with known $Q$ and linear observables $\bm{\gamma}$. 
	Traces are averaged over bins $\Delta$ spanning sub- to supra-kinetic regimes. 
	White and pink noise components $(\epsilon^2,\nu^2)$ follow baseline experimental estimates. 
	Priors on rate constants are log-uniform, and noise parameters use weakly-informative priors.
	
	\subsection*{Comparative benchmarks}
	We implemented three alternatives: (i) midpoint likelihood (single-point per bin), (ii) non-recursive interval likelihood (NR) without boundary conditioning, and (iii) instantaneous-sample approximation. 
	Each method was compared in terms of bias, RMSE, coverage of 95\% credible intervals, and runtime. 
	MacroIR used identical priors and samplers. 
	Model ranking employed Bayesian evidence computed via thermodynamic integration.
	
	\subsection*{Score test}
	Let $S(\theta)=\nabla_\theta \log \mathcal{L}(\theta)$ denote the per-interval score. 
	Under a correct likelihood, $\mathbb{E}[S]=0$ and $\mathrm{Var}(S)=\mathcal{I}(\theta)$ (Fisher information). 
	We estimated $\langle S\rangle$ and $\mathrm{Var}(S)$ across intervals and compared them to plug-in Fisher estimates from MacroIR, evaluating deviations as a measure of informational self-consistency.
	
	\section*{Results}
	
	\subsection*{Synthetic recovery and benchmark comparisons}
	MacroIR recovered unbiased parameters across all $\Delta/\tau_{\min}$ regimes, with calibrated uncertainty and higher evidence for true models. 
	Midpoint and NR approaches displayed increasing bias and under-coverage as $\Delta$ grew, while instantaneous approximations failed when $\Delta\gtrsim\tau_{\min}$.
	
	\subsection*{Score-based self-consistency}
	Only MacroIR yielded mean score $\approx0$ with variance matching Fisher predictions across sampling regimes. 
	Alternative methods showed nonzero mean scores and underestimated information, especially at large $\Delta$, revealing internal inconsistency between their assumed likelihood and the physical averaging process.
	
	\subsection*{Biological case study: P2X receptor kinetics}
	Reanalysis of ultrafast ATP-pulse macroscopic currents from P2X$_2$ receptors showed that MacroIR reproduced experimental responses and correctly ranked candidate kinetic schemes, revealing asymmetric coupling between ligand binding and subunit rotation. 
	These data were previously analyzed using the Macro R approximation \citep{Moffatt2007}; MacroIR reproduces those results in the instantaneous limit and reveals additional asymmetry once interval averaging is incorporated.
	
	\subsection*{Computational performance}
	Computation scales linearly with the number of intervals and as $O(k^3)$ with state dimension, dominated by matrix exponentials and linear solves. 
	Including an explicit analog filter via state-space augmentation increases runtime $\sim$10$\times$ but remains tractable for moderate $k$.
	
	\section*{Discussion}
	
	MacroIR provides a unified Bayesian framework for inference from time-averaged observations governed by continuous-time Markov processes. 
	It generalizes the macroscopic recursive algorithm of Moffatt (2007), which already cast macroscopic current analysis in Bayesian terms, to handle explicit temporal averaging and to compute model evidence. 
	The resulting equations encompass the Kalman recursion as a limiting case but retain their validity beyond linear–Gaussian assumptions.  
	By restoring consistency between the physics of sampling and the statistical model, MacroIR eliminates the structural bias intrinsic to instantaneous-likelihood approaches.
	
	The Bayesian foundation of MacroIR allows immediate extensions to non-Gaussian noise, hierarchical priors, and mixture models, and provides a natural route to Bayesian evidence comparison across kinetic schemes. 
	This generality distinguishes it from purely algorithmic filters and repositions the theory of macroscopic inference within the domain of full probabilistic reasoning.
	
	\subsection*{Limitations and outlook}
	(1) Scaling may become demanding for very large $k$; sparsity or exponential-Krylov methods could improve performance. 
	(2) Explicit analog filtering requires state-space augmentation or kernel convolution and increases runtime. 
	(3) Robustness to model misspecification and application to hierarchical or mixture CTMCs remain open avenues. 
	(4) Nonlinear observables will require moment-closure or variational approximations beyond the Gaussian level.
	
	\section*{Data and Code Availability}
	Code and scripts to reproduce all results will be released publicly (GitHub/Zenodo) at submission, including minimal working examples and figure-generation notebooks.
	
	\section*{Acknowledgements}
	I thank colleagues for discussions on Bayesian conditioning of time-averaged processes and the INQUIMAE computing facility for support.
	
	\section*{Author Contributions}
	LM conceived the study, developed the theory, implemented algorithms, analyzed data, and wrote the manuscript.
	
	\bibliography{biblio}
	
\end{document}
